\clearpage
\section*{附件}
本文采用本届A题题面、技术附录与随题官方评估程序。随文给出结果文件、关键实现和逐图附表；100图逐例工期、额外搬运、Cache配对及参数扰动结果列于后续附录。
\begin{table}[H]\centering\small\auxtablecaption{结果与文件索引}\label{tab:support}
\setlength{\tabcolsep}{4pt}
\begin{tabularx}{\linewidth}{l>{\hsize=1.25\hsize}Y>{\hsize=.75\hsize}Y}\toprule
类别 & 相对路径或文件名 & 内容\\\midrule
主矩阵 & \path{data/v7_fast_20260927/main_results.csv} & 100图、三问、1至5核最终与初始结果\\
候选比较 & \path{data/v7_fast_20260927/candidate_effects.csv} & 1500组初始/最终来源与降幅\\
Cache配对 & \path{data/v7_fast_20260927/cache_pairs.csv} & 500组B最终计划起始的三段工期比及$C(X_B)$/$C(X_C)$分离命中字段\\
参数实验 & \path{data/v7_fast_20260927/sensitivity_summary.csv} & 六图四变体的固定与重优化结果\\
预算实验 & \path{budget_sensitivity_20260927_v2/budget_aggregate.csv} & 六张代表图在2、5、10、20次评价预算下的均值、成功率与累计耗时\\
失败记录 & \path{data/v7_fast_20260927/sensitivity_failures.csv} & v7候选失败记录（本轮无失败候选）\\
关键实现 & \path{support/experiment/v7_fast_20260927/src/huawei_code/} & 图预处理、候选构造、评分和三问接口\\
官方程序 & \path{support/experiment/v7_fast_20260927/vendor/official_evaluator_fast/} & 题目提供的事件评价程序\\
结果汇总 & \path{data/v7_fast_20260927/report_manifest.json} & 组合覆盖与字段汇总\\
v7汇总 & \path{data/v7_fast_20260927/paper_results.md} & final\_selection主矩阵与均值统计\\\bottomrule
\end{tabularx}\end{table}
逐图附表报告最终工期、相对整图单核 A0 速度比和额外搬运。\path{data/v7_fast_20260927/main_results.csv} 另外保存 \texttt{baseline\_makespan}、\texttt{initial\_makespan}、\texttt{final\_makespan}、候选来源及方案哈希；\texttt{final\_partition\_added\_copy\_bytes} 与 \texttt{final\_spill\_added\_copy\_bytes} 分别给出切分搬运和 Spill 搬运。\path{data/v7_fast_20260927/candidate_effects.csv} 保存首次成功候选到最终候选的工期与相对变化，附表字段与结果文件的范围由此分开。
\clearpage
\appendix
\ctexset{section={aftername={\quad}}}
\section{关键源代码节选}\label{app:source-code}
以下片段对应本文的四个模型步骤：计算图预处理、候选方案构造、轻量排序和评价接口，仅列出能够对应公式、约束和事件定义的关键实现。

\begingroup
\setmonofont{Menlo}
\fvset{fontsize=\scriptsize,baselinestretch=1,breaklines=true,breakanywhere=true,breakanywheresymbolpre={},breaksymbolleft={},numbers=left,numbersep=4pt,xleftmargin=2em}

\subsection{计算图与任务优先级}
\noindent\textbf{计算图与任务优先级：}张量依赖收缩、关键路径和释放优先级。\par\nobreak
\VerbatimInput[firstline=22,lastline=70,firstnumber=22]{support/experiment/v7_fast_20260927/src/huawei_code/graph.py}
\VerbatimInput[firstline=91,lastline=107,firstnumber=91]{support/experiment/v7_fast_20260927/src/huawei_code/graph.py}
\VerbatimInput[firstline=110,lastline=141,firstnumber=110]{support/experiment/v7_fast_20260927/src/huawei_code/graph.py}
\VerbatimInput[firstline=144,lastline=170,firstnumber=144]{support/experiment/v7_fast_20260927/src/huawei_code/graph.py}

\subsection{候选方案构造与约束}
\noindent\textbf{候选方案构造与约束：}计划规范化、场景化依赖合法性检查、流水线切块和锚点方案。\par\nobreak
\VerbatimInput[firstline=19,lastline=96,firstnumber=19]{support/experiment/v7_fast_20260927/src/huawei_code/plan.py}
\VerbatimInput[firstline=110,lastline=150,firstnumber=110]{support/experiment/v7_fast_20260927/src/huawei_code/plan.py}

\noindent\textbf{候选变换：}局部移动、聚合、重排及共享张量的Cache候选。\par\nobreak
\VerbatimInput[firstline=13,lastline=24,firstnumber=13]{support/experiment/v7_fast_20260927/src/huawei_code/candidates.py}
\VerbatimInput[firstline=27,lastline=78,firstnumber=27]{support/experiment/v7_fast_20260927/src/huawei_code/candidates.py}
\VerbatimInput[firstline=81,lastline=97,firstnumber=81]{support/experiment/v7_fast_20260927/src/huawei_code/candidates.py}
\VerbatimInput[firstline=100,lastline=140,firstnumber=100]{support/experiment/v7_fast_20260927/src/huawei_code/candidates.py}

\subsection{轻量排序与评估程序评价}
\noindent\textbf{轻量排序：}互斥结构搬运、容量压力、逻辑COPY查询代理及候选排序。\par\nobreak
\VerbatimInput[firstline=31,lastline=80,firstnumber=31]{support/experiment/v7_fast_20260927/src/huawei_code/scoring.py}
\VerbatimInput[firstline=83,lastline=119,firstnumber=83]{support/experiment/v7_fast_20260927/src/huawei_code/scoring.py}
\VerbatimInput[firstline=127,lastline=206,firstnumber=127]{support/experiment/v7_fast_20260927/src/huawei_code/scoring.py}

\noindent\textbf{评价接口：}调用题面评价程序并分别返回A、B、C三问的工期、搬运和命中字段。\par\nobreak
\VerbatimInput[firstline=11,lastline=16,firstnumber=11]{support/experiment/v7_fast_20260927/src/huawei_code/official.py}
\VerbatimInput[firstline=44,lastline=58,firstnumber=44]{support/experiment/v7_fast_20260927/src/huawei_code/official.py}
\endgroup

\section{数据与指标说明}\label{app:provenance}
本文使用题面给出的100张计算图、硬件配置和评价程序。主表包含100图、1至5核和A、B、C三种场景的1500条记录；其中A、B、C各核数均为100条；另有500组B/C同计划配对和24组Cache参数固定计划评价。表中工期、搬运和命中率均由相应评价结果汇总。

图中区间以100图有放回重采样1500次所得经验分位数表示实例组成对均值的影响。加速比均由原始工期计算。统一基准工期为整图单Task单核 A0 的 $T_{1,i}$；A、B 单核曲线按题面显示为 1.000，逐例表的 A/B 单核单元仍列各自最终已评计划的实际工期，C 单核单元列实际工期比。以\Case{053}为例，A/B/C 的单核实际工期分别为 1471202、1448004、1426552 cycle，统一基准为 1472686 cycle。

\section{三问逐图结果}\label{app:results}
表\ref{tab:detail-1}至表\ref{tab:detail-5}共列出100图、5种核数、3问的1500条最终记录。
每个$T_r/S_r$单元中的斜线分别分隔工期（cycle）和相对整图单核 A0 的加速比；$D_r$是官方额外搬运字节。A、B、C单核均保留主矩阵实际工期比；题面展示时可按题面要求单独归一化；三种场景的单核实际工期均来自对应最终已评计划。
逐例表只列出最终工期、相对整图单核 A0 的加速比和额外数据搬运量。初始工期、候选来源、方案哈希、切分搬运与 Spill 搬运见 \path{data/v7_fast_20260927/main_results.csv}；\path{data/v7_fast_20260927/candidate_effects.csv} 另列首个成功候选至最终候选的逐组工期变化。
\begingroup
\small
\setlength{\tabcolsep}{2.5pt}
\setlength{\LTleft}{\fill}\setlength{\LTright}{\fill}
\renewcommand{\arraystretch}{1.03}
\begin{longtable}{crrrrrr}
\caption{1核三问逐例工期/加速比与额外搬运（100图）}\label{tab:detail-1}\\
\toprule
图 & $T_A/S_A$ & $D_A$/B & $T_B/S_B$ & $D_B$/B & $T_C/S_C$ & $D_C$/B \\
\midrule
\endfirsthead
\multicolumn{7}{c}{续表}\\
\toprule
图 & $T_A/S_A$ & $D_A$/B & $T_B/S_B$ & $D_B$/B & $T_C/S_C$ & $D_C$/B \\
\midrule
\endhead
\midrule\multicolumn{7}{r}{续下页}\\
\endfoot
\bottomrule
\endlastfoot
001 & 233110/1.000 & 0 & 233110/1.000 & 0 & 233110/1.000 & 0 \\
002 & 261945/1.000 & 0 & 260509/1.000 & 2311680 & 254213/1.030 & 2311680 \\
003 & 587919/1.000 & 649106 & 592981/1.000 & 0 & 591743/1.018 & 0 \\
004 & 392357/1.000 & 0 & 392195/1.000 & 2049024 & 387184/1.013 & 2049024 \\
005 & 95618/1.000 & 0 & 93310/1.000 & 0 & 93171/1.026 & 0 \\
006 & 80821/1.000 & 1229146 & 73949/1.000 & 278560 & 73949/1.120 & 278560 \\
007 & 985964/1.000 & 9369600 & 1012562/1.000 & 0 & 1012562/1.000 & 0 \\
008 & 413595/1.000 & 11596800 & 269325/1.000 & 3760128 & 269325/1.810 & 3760128 \\
009 & 153261/1.000 & 0 & 153201/1.000 & 0 & 153201/1.000 & 0 \\
010 & 78330/1.000 & 1385412 & 81326/1.000 & 224256 & 79774/1.090 & 224256 \\
011 & 184697/1.000 & 0 & 184697/1.000 & 0 & 184697/1.000 & 0 \\
012 & 58890/1.000 & 3456 & 42251/1.000 & 8192 & 42239/1.399 & 8192 \\
013 & 263949/1.000 & 0 & 263949/1.000 & 0 & 263949/1.000 & 0 \\
014 & 17698038/1.000 & 71385504 & 17698626/1.000 & 71010720 & 17689602/1.001 & 71010720 \\
015 & 172300/1.000 & 1778432 & 167767/1.000 & 507940 & 167630/1.058 & 507940 \\
016 & 7715523/1.000 & 0 & 7715523/1.000 & 0 & 7715523/1.000 & 0 \\
017 & 162810/1.000 & 22528 & 149759/1.000 & 180224 & 149367/1.090 & 180224 \\
018 & 1004367/1.000 & 1970176 & 990647/1.000 & 5537792 & 990647/1.014 & 5537792 \\
019 & 75090/1.000 & 1050880 & 70264/1.000 & 6144 & 70264/1.155 & 6144 \\
020 & 532746/1.000 & 0 & 532746/1.000 & 0 & 532746/1.000 & 0 \\
021 & 6520170/1.000 & 4702208 & 6567937/1.000 & 4501504 & 6558837/1.001 & 4501504 \\
022 & 56780/1.000 & 0 & 54723/1.000 & 88064 & 54591/1.040 & 78848 \\
023 & 110425/1.000 & 897024 & 99763/1.000 & 75776 & 93197/1.316 & 75776 \\
024 & 2555343/1.000 & 0 & 2555343/1.000 & 0 & 2555343/1.000 & 0 \\
025 & 4828832/1.000 & 44894208 & 4863026/1.000 & 0 & 4863026/1.000 & 0 \\
026 & 97011/1.000 & 49184 & 95759/1.000 & 942080 & 88492/1.121 & 1064960 \\
027 & 2325280/1.000 & 1425408 & 2304836/1.000 & 3452928 & 2304836/1.026 & 3452928 \\
028 & 43242729/1.000 & 206054016 & 42901220/1.000 & 161588040 & 41700132/1.051 & 161588040 \\
029 & 37555/1.000 & 0 & 37555/1.000 & 0 & 37555/1.000 & 0 \\
030 & 2808034/1.000 & 5519360 & 2808034/1.000 & 5519360 & 2806776/1.000 & 5395456 \\
031 & 1410105/1.000 & 3829248 & 1414369/1.000 & 3491328 & 1414159/1.000 & 3491328 \\
032 & 41443/1.000 & 0 & 34141/1.000 & 10256 & 34141/1.214 & 0 \\
033 & 1774196/1.000 & 4592640 & 1774196/1.000 & 4592640 & 1774080/1.000 & 4592640 \\
034 & 839876/1.000 & 390672 & 842996/1.000 & 1079832 & 841851/1.001 & 746624 \\
035 & 169571/1.000 & 0 & 169168/1.000 & 0 & 168511/1.006 & 0 \\
036 & 931510/1.000 & 0 & 931510/1.000 & 0 & 931510/1.000 & 0 \\
037 & 210001/1.000 & 0 & 210001/1.000 & 0 & 210001/1.000 & 0 \\
038 & 1363201/1.000 & 4233216 & 1364555/1.000 & 5627904 & 1359544/1.004 & 4697600 \\
039 & 1430004/1.000 & 7954432 & 1435262/1.000 & 7929856 & 1426335/1.006 & 8175616 \\
040 & 253454/1.000 & 0 & 252193/1.000 & 137892 & 249653/1.015 & 137892 \\
041 & 6003957/1.000 & 30898176 & 6003957/1.000 & 30898176 & 6001575/1.000 & 30898176 \\
042 & 1686740/1.000 & 6629376 & 1686740/1.000 & 6629376 & 1675633/1.007 & 6629376 \\
043 & 677656/1.000 & 49160 & 657411/1.000 & 530098 & 656100/1.033 & 530098 \\
044 & 96408/1.000 & 109824 & 88538/1.000 & 0 & 88538/1.744 & 0 \\
045 & 125614/1.000 & 0 & 125614/1.000 & 0 & 125614/1.000 & 0 \\
046 & 233586/1.000 & 319488 & 226233/1.000 & 163840 & 226233/1.222 & 163840 \\
047 & 324416/1.000 & 257106 & 323597/1.000 & 0 & 322738/1.012 & 0 \\
048 & 222220/1.000 & 0 & 222220/1.000 & 0 & 222220/1.000 & 0 \\
049 & 196276/1.000 & 0 & 195430/1.000 & 41984 & 194548/1.009 & 41984 \\
050 & 149690/1.000 & 0 & 149242/1.000 & 409608 & 148876/1.005 & 836812 \\
051 & 607628/1.000 & 0 & 607628/1.000 & 0 & 607628/1.000 & 0 \\
052 & 179776/1.000 & 1152896 & 159895/1.000 & 70848 & 158577/1.137 & 70848 \\
053 & 1415374/1.000 & 1759362 & 1448004/1.000 & 4047952 & 1421619/1.036 & 4042192 \\
054 & 815051/1.000 & 0 & 796718/1.000 & 165700 & 792519/1.028 & 213576 \\
055 & 118835/1.000 & 69632 & 118835/1.000 & 69632 & 114628/1.037 & 57396 \\
056 & 277285/1.000 & 0 & 277285/1.000 & 0 & 277285/1.000 & 0 \\
057 & 59726/1.000 & 94216 & 59247/1.000 & 200740 & 59096/1.011 & 274492 \\
058 & 4527949/1.000 & 28782592 & 4539717/1.000 & 28696576 & 4533133/1.001 & 28721152 \\
059 & 818974/1.000 & 2072576 & 829473/1.000 & 1691648 & 815938/1.017 & 1691648 \\
060 & 886258/1.000 & 2398344 & 885660/1.000 & 1869960 & 883441/1.006 & 1977582 \\
061 & 77191/1.000 & 0 & 77191/1.000 & 0 & 77191/1.000 & 0 \\
062 & 2821994/1.000 & 42112512 & 2854809/1.000 & 0 & 2854809/1.000 & 0 \\
063 & 1070397/1.000 & 0 & 1070397/1.000 & 0 & 1070397/1.000 & 0 \\
064 & 19034/1.000 & 12198 & 19116/1.000 & 0 & 19116/1.000 & 0 \\
065 & 51959/1.000 & 8200 & 39901/1.000 & 0 & 39824/1.333 & 0 \\
066 & 359070/1.000 & 0 & 357763/1.000 & 0 & 356326/1.008 & 0 \\
067 & 63834834/1.000 & 264449408 & 63397576/1.000 & 256305152 & 63397576/1.007 & 256305152 \\
068 & 362095/1.000 & 0 & 360298/1.000 & 0 & 359166/1.008 & 0 \\
069 & 24642/1.000 & 61828 & 22905/1.000 & 0 & 22655/1.110 & 0 \\
070 & 115306/1.000 & 73728 & 115088/1.000 & 0 & 104334/1.124 & 0 \\
071 & 18919/1.000 & 0 & 17557/1.000 & 0 & 17401/1.087 & 0 \\
072 & 23897828/1.000 & 108613680 & 23897828/1.000 & 108613680 & 23890756/1.000 & 108613680 \\
073 & 9876885/1.000 & 136869376 & 9870212/1.000 & 136743424 & 9870212/1.021 & 136743424 \\
074 & 1682239/1.000 & 6070280 & 1682239/1.000 & 6070280 & 1678120/1.002 & 6070280 \\
075 & 1441809/1.000 & 0 & 1433562/1.000 & 0 & 1427745/1.010 & 0 \\
076 & 6180531/1.000 & 25608192 & 6185113/1.000 & 28280832 & 6182761/1.000 & 26796032 \\
077 & 495320/1.000 & 0 & 495320/1.000 & 0 & 494044/1.003 & 0 \\
078 & 2249980/1.000 & 9199104 & 2232856/1.000 & 8879232 & 2142955/1.071 & 4056192 \\
079 & 30952656/1.000 & 32555008 & 31052251/1.000 & 32448512 & 31052251/1.000 & 32448512 \\
080 & 292086/1.000 & 667648 & 292138/1.000 & 548864 & 291826/1.001 & 548864 \\
081 & 151795/1.000 & 0 & 151795/1.000 & 0 & 151795/1.000 & 0 \\
082 & 716933/1.000 & 148102 & 715588/1.000 & 0 & 709534/1.016 & 0 \\
083 & 1057993/1.000 & 13934688 & 876669/1.000 & 916480 & 875735/1.208 & 839680 \\
084 & 1979969/1.000 & 29566656 & 1371156/1.000 & 9676800 & 1371156/1.829 & 9676800 \\
085 & 3187899/1.000 & 0 & 3177963/1.000 & 355712 & 3173461/1.005 & 355712 \\
086 & 118296/1.000 & 84966 & 115732/1.000 & 0 & 114476/1.034 & 0 \\
087 & 3974071/1.000 & 7902720 & 3974071/1.000 & 7902720 & 3969779/1.001 & 7902720 \\
088 & 243623/1.000 & 0 & 243542/1.000 & 0 & 242071/1.006 & 0 \\
089 & 693092/1.000 & 0 & 689572/1.000 & 0 & 689572/1.005 & 0 \\
090 & 1092711/1.000 & 15390192 & 1087442/1.000 & 14905728 & 919256/1.282 & 6461952 \\
091 & 9186769/1.000 & 53947008 & 9186769/1.000 & 53947008 & 9166672/1.002 & 48764544 \\
092 & 4070807/1.000 & 8853504 & 4010291/1.000 & 5283840 & 4004200/1.203 & 4695040 \\
093 & 74205/1.000 & 0 & 74205/1.000 & 0 & 74205/1.000 & 0 \\
094 & 94892/1.000 & 0 & 91195/1.000 & 122880 & 87488/1.085 & 122880 \\
095 & 2084853/1.000 & 0 & 2070244/1.000 & 49848192 & 1161033/1.796 & 16551936 \\
096 & 84194/1.000 & 0 & 77108/1.000 & 0 & 75108/1.121 & 0 \\
097 & 11491509/1.000 & 10584064 & 11491509/1.000 & 10584064 & 11491509/1.000 & 10584064 \\
098 & 177430/1.000 & 0 & 176752/1.000 & 364588 & 174782/1.015 & 385068 \\
099 & 1619135/1.000 & 5885952 & 1622999/1.000 & 5484544 & 1618902/1.003 & 5484544 \\
100 & 148339/1.000 & 5640 & 148885/1.000 & 0 & 142753/1.059 & 0 \\
\end{longtable}

\begin{longtable}{crrrrrr}
\caption{2核三问逐例工期/加速比与额外搬运（100图）}\label{tab:detail-2}\\
\toprule
图 & $T_A/S_A$ & $D_A$/B & $T_B/S_B$ & $D_B$/B & $T_C/S_C$ & $D_C$/B \\
\midrule
\endfirsthead
\multicolumn{7}{c}{续表}\\
\toprule
图 & $T_A/S_A$ & $D_A$/B & $T_B/S_B$ & $D_B$/B & $T_C/S_C$ & $D_C$/B \\
\midrule
\endhead
\midrule\multicolumn{7}{r}{续下页}\\
\endfoot
\bottomrule
\endlastfoot
001 & 116868/1.995 & 1152 & 116868/1.995 & 1152 & 116868/1.995 & 1152 \\
002 & 230910/1.134 & 5776896 & 165091/1.587 & 3614208 & 164754/1.590 & 3614208 \\
003 & 538381/1.119 & 98304 & 538381/1.119 & 98304 & 531381/1.133 & 135302 \\
004 & 199074/1.971 & 110592 & 199074/1.971 & 110592 & 199074/1.971 & 110592 \\
005 & 95618/1.000 & 0 & 93310/1.025 & 0 & 93060/1.027 & 170916 \\
006 & 42696/1.940 & 56200 & 42696/1.940 & 56200 & 41180/2.012 & 56200 \\
007 & 507494/1.995 & 27648 & 507494/1.995 & 27648 & 507494/1.995 & 27648 \\
008 & 244266/1.996 & 0 & 244266/1.996 & 0 & 244266/1.996 & 0 \\
009 & 92198/1.662 & 0 & 92198/1.662 & 0 & 92198/1.662 & 0 \\
010 & 45272/1.921 & 48384 & 45272/1.921 & 48384 & 41774/2.082 & 48384 \\
011 & 92646/1.994 & 180224 & 92646/1.994 & 180224 & 92646/1.994 & 180224 \\
012 & 30058/1.966 & 29376 & 30058/1.966 & 29376 & 29156/2.027 & 29376 \\
013 & 132198/1.997 & 8192 & 132198/1.997 & 8192 & 132198/1.997 & 8192 \\
014 & 8897730/1.989 & 71757240 & 8897730/1.989 & 71757240 & 8888616/1.991 & 71757240 \\
015 & 90509/1.960 & 199948 & 90509/1.960 & 199948 & 87390/2.030 & 199948 \\
016 & 7715523/1.000 & 0 & 7712915/1.000 & 2162724 & 7712915/1.000 & 2162724 \\
017 & 81772/1.991 & 87552 & 81772/1.991 & 87552 & 81772/1.991 & 87552 \\
018 & 504468/1.991 & 3253248 & 504468/1.991 & 3253248 & 502571/1.998 & 3253248 \\
019 & 41382/1.961 & 27168 & 41382/1.961 & 27168 & 39509/2.054 & 27168 \\
020 & 266580/1.998 & 0 & 266580/1.998 & 0 & 266580/1.998 & 0 \\
021 & 3194570/2.056 & 3174400 & 3194570/2.056 & 3174400 & 3193634/2.057 & 3174400 \\
022 & 28324/2.005 & 6912 & 28324/2.005 & 6912 & 28324/2.005 & 6912 \\
023 & 61177/2.004 & 84096 & 61177/2.004 & 84096 & 60171/2.038 & 84096 \\
024 & 2555343/1.000 & 0 & 2552735/1.001 & 2162724 & 2552735/1.001 & 2162724 \\
025 & 2431840/2.000 & 27648 & 2431840/2.000 & 27648 & 2431840/2.000 & 27648 \\
026 & 55057/1.802 & 21936 & 55057/1.802 & 21936 & 51429/1.929 & 21936 \\
027 & 1116984/2.116 & 2211840 & 1120962/2.109 & 1744896 & 1120962/2.109 & 1744896 \\
028 & 23638628/1.854 & 277473600 & 23638628/1.854 & 277473600 & 23615145/1.856 & 277473600 \\
029 & 19740/1.902 & 0 & 19740/1.902 & 0 & 19740/1.902 & 0 \\
030 & 1399149/2.007 & 6877184 & 1399149/2.007 & 6877184 & 1398912/2.007 & 6877184 \\
031 & 712596/1.985 & 5250048 & 712596/1.985 & 5250048 & 706497/2.002 & 4666368 \\
032 & 21040/1.970 & 6216 & 21040/1.970 & 6216 & 20711/2.001 & 6216 \\
033 & 889697/1.994 & 6664224 & 889697/1.994 & 6664224 & 887951/1.998 & 6664224 \\
034 & 418129/2.016 & 1648128 & 418129/2.016 & 1648128 & 417777/2.018 & 1648128 \\
035 & 95406/1.777 & 0 & 95406/1.777 & 0 & 95406/1.777 & 0 \\
036 & 466068/1.999 & 1152 & 466068/1.999 & 1152 & 466068/1.999 & 1152 \\
037 & 105374/1.993 & 204800 & 105374/1.993 & 204800 & 105374/1.993 & 204800 \\
038 & 674500/2.023 & 6790656 & 674500/2.023 & 6790656 & 674381/2.023 & 6173696 \\
039 & 708561/2.026 & 11067392 & 708561/2.026 & 11067392 & 703954/2.039 & 11186176 \\
040 & 137746/1.840 & 53760 & 137746/1.840 & 53760 & 137487/1.843 & 57856 \\
041 & 2983392/2.012 & 30750720 & 2983392/2.012 & 30750720 & 2981318/2.014 & 30750720 \\
042 & 845274/1.995 & 6540576 & 844150/1.998 & 6560160 & 843098/2.001 & 6570528 \\
043 & 347613/1.949 & 417800 & 347613/1.949 & 417800 & 347613/1.949 & 417800 \\
044 & 96408/1.602 & 109824 & 68899/2.241 & 120384 & 67862/2.275 & 128576 \\
045 & 63059/1.992 & 0 & 63059/1.992 & 0 & 63052/1.992 & 12288 \\
046 & 160984/1.717 & 4513952 & 166578/1.660 & 3521632 & 125148/2.209 & 2316128 \\
047 & 324416/1.006 & 257106 & 321355/1.016 & 1146116 & 321355/1.016 & 1146116 \\
048 & 222220/1.000 & 0 & 219890/1.011 & 330436 & 219037/1.015 & 294724 \\
049 & 160206/1.225 & 0 & 160206/1.225 & 0 & 160206/1.225 & 0 \\
050 & 108121/1.384 & 3456 & 108121/1.384 & 3456 & 107981/1.386 & 11652 \\
051 & 607628/1.000 & 0 & 607628/1.000 & 0 & 607628/1.000 & 0 \\
052 & 90903/1.983 & 102912 & 90903/1.983 & 102912 & 90903/1.983 & 102912 \\
053 & 774929/1.900 & 5959552 & 774929/1.900 & 5959552 & 759346/1.939 & 5308480 \\
054 & 517960/1.574 & 49152 & 517960/1.574 & 49152 & 513641/1.587 & 49152 \\
055 & 59839/1.986 & 41472 & 59839/1.986 & 41472 & 59839/1.986 & 41472 \\
056 & 265710/1.044 & 0 & 265710/1.044 & 0 & 265710/1.044 & 0 \\
057 & 30183/1.979 & 22304 & 30183/1.979 & 22304 & 29658/2.014 & 22304 \\
058 & 2271316/1.999 & 33431552 & 2268376/2.001 & 33435648 & 2266610/2.003 & 33435648 \\
059 & 369318/2.246 & 360448 & 369318/2.246 & 360448 & 369318/2.246 & 360448 \\
060 & 443955/2.002 & 2603148 & 443955/2.002 & 2603148 & 437319/2.032 & 2715852 \\
061 & 39298/1.964 & 34560 & 39298/1.964 & 34560 & 38792/1.990 & 46848 \\
062 & 2821994/1.012 & 42112512 & 2521494/1.132 & 68806656 & 2387918/1.196 & 68368896 \\
063 & 795582/1.345 & 2116608 & 909394/1.177 & 23271936 & 874911/1.223 & 24056832 \\
064 & 19034/1.004 & 12198 & 19116/1.000 & 0 & 19116/1.000 & 0 \\
065 & 26610/1.994 & 13968 & 26610/1.994 & 13968 & 26551/1.999 & 13968 \\
066 & 248146/1.447 & 0 & 248146/1.447 & 0 & 247485/1.451 & 0 \\
067 & 33768704/1.890 & 264024576 & 33768704/1.890 & 264024576 & 32164982/1.985 & 260124928 \\
068 & 346246/1.046 & 0 & 303697/1.192 & 1343990 & 302692/1.196 & 1343990 \\
069 & 24642/1.020 & 61828 & 22905/1.098 & 0 & 22655/1.110 & 0 \\
070 & 60121/1.951 & 129536 & 60121/1.951 & 129536 & 59140/1.983 & 129536 \\
071 & 18919/1.000 & 0 & 17557/1.078 & 0 & 17401/1.087 & 0 \\
072 & 11833111/2.020 & 113945376 & 11833111/2.020 & 113945376 & 11819724/2.022 & 113945376 \\
073 & 5980423/1.686 & 141121024 & 5980423/1.686 & 141121024 & 5176545/1.948 & 138831104 \\
074 & 821757/2.047 & 7077888 & 821757/2.047 & 7077888 & 818681/2.055 & 7086080 \\
075 & 1441809/1.000 & 0 & 1397410/1.032 & 2016420 & 1389088/1.038 & 2016420 \\
076 & 3065214/2.018 & 31095808 & 3065214/2.018 & 31095808 & 3063504/2.019 & 31105024 \\
077 & 248734/1.991 & 57344 & 248734/1.991 & 57344 & 248734/1.991 & 57344 \\
078 & 1177840/1.949 & 13146912 & 1260358/1.822 & 12535344 & 1149984/1.997 & 11180376 \\
079 & 15446486/2.010 & 31780864 & 15446486/2.010 & 31780864 & 15446486/2.010 & 31780864 \\
080 & 145926/2.002 & 2191360 & 145926/2.002 & 2191360 & 145674/2.005 & 2191360 \\
081 & 76860/1.975 & 0 & 76860/1.975 & 0 & 76860/1.975 & 0 \\
082 & 713527/1.011 & 376592 & 708711/1.017 & 234630 & 708711/1.017 & 234630 \\
083 & 644672/1.641 & 14241216 & 644672/1.641 & 14241216 & 465958/2.271 & 7847648 \\
084 & 1255052/1.998 & 0 & 1255052/1.998 & 0 & 1255052/1.998 & 0 \\
085 & 3187899/1.000 & 0 & 3168404/1.006 & 6815178 & 3157230/1.010 & 6815178 \\
086 & 118296/1.001 & 84966 & 115732/1.023 & 0 & 114476/1.034 & 0 \\
087 & 1985234/2.002 & 9999360 & 1984436/2.003 & 10008576 & 1982828/2.004 & 10008576 \\
088 & 221995/1.097 & 0 & 203492/1.197 & 1597604 & 203492/1.197 & 1597604 \\
089 & 346832/1.998 & 238720 & 346832/1.998 & 238720 & 346763/1.999 & 238720 \\
090 & 769224/1.532 & 21498480 & 677936/1.738 & 11583648 & 549431/2.145 & 5289120 \\
091 & 4571759/2.009 & 57904096 & 4571759/2.009 & 57904096 & 4569395/2.011 & 57904096 \\
092 & 2757789/1.746 & 67180384 & 2828252/1.703 & 66389088 & 2358868/2.041 & 66389088 \\
093 & 37827/1.962 & 8192 & 37827/1.962 & 8192 & 37827/1.962 & 8192 \\
094 & 58760/1.615 & 106272 & 48073/1.974 & 106272 & 46927/2.022 & 106272 \\
095 & 1042890/1.999 & 0 & 1042890/1.999 & 0 & 1042890/1.999 & 0 \\
096 & 44010/1.913 & 67392 & 44010/1.913 & 67392 & 41802/2.014 & 67392 \\
097 & 5617602/2.046 & 11796480 & 5617602/2.046 & 11796480 & 5617602/2.046 & 11796480 \\
098 & 89756/1.977 & 0 & 89756/1.977 & 0 & 89756/1.977 & 0 \\
099 & 808916/2.006 & 7254016 & 808916/2.006 & 7254016 & 806053/2.014 & 7254016 \\
100 & 81080/1.865 & 16896 & 81080/1.865 & 16896 & 77421/1.953 & 16896 \\
\end{longtable}

\begin{longtable}{crrrrrr}
\caption{3核三问逐例工期/加速比与额外搬运（100图）}\label{tab:detail-3}\\
\toprule
图 & $T_A/S_A$ & $D_A$/B & $T_B/S_B$ & $D_B$/B & $T_C/S_C$ & $D_C$/B \\
\midrule
\endfirsthead
\multicolumn{7}{c}{续表}\\
\toprule
图 & $T_A/S_A$ & $D_A$/B & $T_B/S_B$ & $D_B$/B & $T_C/S_C$ & $D_C$/B \\
\midrule
\endhead
\midrule\multicolumn{7}{r}{续下页}\\
\endfoot
\bottomrule
\endlastfoot
001 & 78545/2.968 & 2304 & 78545/2.968 & 2304 & 78545/2.968 & 2304 \\
002 & 173068/1.514 & 594432 & 131349/1.994 & 4211712 & 127539/2.054 & 4294656 \\
003 & 538381/1.119 & 98304 & 524210/1.149 & 3647050 & 522594/1.152 & 3579466 \\
004 & 134043/2.927 & 221184 & 134043/2.927 & 221184 & 134043/2.927 & 221184 \\
005 & 95618/1.000 & 0 & 88696/1.078 & 1492040 & 86911/1.100 & 1492040 \\
006 & 31152/2.659 & 30080 & 31152/2.659 & 30080 & 31152/2.659 & 30080 \\
007 & 338094/2.995 & 55296 & 338094/2.995 & 55296 & 338094/2.995 & 55296 \\
008 & 163359/2.985 & 0 & 163359/2.985 & 0 & 163359/2.985 & 0 \\
009 & 92198/1.662 & 0 & 91278/1.679 & 9254 & 91278/1.679 & 9254 \\
010 & 32786/2.653 & 38016 & 32786/2.653 & 38016 & 30412/2.860 & 91008 \\
011 & 69625/2.653 & 1081344 & 69279/2.666 & 1095680 & 68637/2.691 & 1097728 \\
012 & 20661/2.860 & 39168 & 20661/2.860 & 39168 & 20561/2.874 & 41472 \\
013 & 88351/2.988 & 16384 & 88351/2.988 & 16384 & 88351/2.988 & 16384 \\
014 & 5980663/2.959 & 77770920 & 5980663/2.959 & 77770920 & 5942569/2.978 & 77770920 \\
015 & 60578/2.928 & 290060 & 60578/2.928 & 290060 & 60578/2.928 & 290060 \\
016 & 7715523/1.000 & 0 & 7712915/1.000 & 2162724 & 7712915/1.000 & 2162724 \\
017 & 55389/2.939 & 175104 & 55389/2.939 & 175104 & 55154/2.952 & 178176 \\
018 & 341047/2.945 & 4093952 & 341047/2.945 & 4093952 & 338488/2.967 & 4093952 \\
019 & 28657/2.832 & 54336 & 28657/2.832 & 54336 & 26572/3.054 & 54336 \\
020 & 177950/2.994 & 0 & 177950/2.994 & 0 & 177950/2.994 & 0 \\
021 & 2144024/3.063 & 7766016 & 2144024/3.063 & 7766016 & 2142612/3.065 & 8163328 \\
022 & 19026/2.984 & 13824 & 19026/2.984 & 13824 & 18965/2.994 & 13824 \\
023 & 40897/2.998 & 207552 & 40897/2.998 & 207552 & 40897/2.998 & 207552 \\
024 & 2555343/1.000 & 0 & 2552735/1.001 & 2162724 & 2552735/1.001 & 2162724 \\
025 & 1622164/2.998 & 55296 & 1622164/2.998 & 55296 & 1622138/2.998 & 62976 \\
026 & 41665/2.381 & 95664 & 41665/2.381 & 95664 & 41665/2.381 & 95664 \\
027 & 837049/2.824 & 3686400 & 831811/2.842 & 4139008 & 831365/2.843 & 4139008 \\
028 & 16525969/2.653 & 277797312 & 16770711/2.614 & 277369056 & 16770711/2.614 & 277369056 \\
029 & 12743/2.947 & 0 & 12743/2.947 & 0 & 12743/2.947 & 0 \\
030 & 933857/3.007 & 9242624 & 933857/3.007 & 9242624 & 932731/3.011 & 9250816 \\
031 & 473822/2.985 & 7544832 & 473822/2.985 & 7544832 & 473000/2.990 & 7544832 \\
032 & 14175/2.924 & 12432 & 14175/2.924 & 12432 & 14175/2.924 & 12432 \\
033 & 592065/2.997 & 6990912 & 592065/2.997 & 6990912 & 590152/3.006 & 6990912 \\
034 & 276994/3.043 & 2830336 & 276994/3.043 & 2830336 & 276560/3.048 & 2830336 \\
035 & 95406/1.777 & 0 & 95406/1.777 & 0 & 95406/1.777 & 0 \\
036 & 311345/2.992 & 2304 & 311345/2.992 & 2304 & 311345/2.992 & 2304 \\
037 & 74038/2.836 & 1064960 & 74038/2.836 & 1064960 & 74038/2.836 & 1064960 \\
038 & 452248/3.017 & 8097792 & 452248/3.017 & 8097792 & 452022/3.019 & 8097792 \\
039 & 485565/2.956 & 14323712 & 485565/2.956 & 14323712 & 482066/2.977 & 14331904 \\
040 & 137746/1.840 & 53760 & 137746/1.840 & 53760 & 137487/1.843 & 57856 \\
041 & 1991911/3.014 & 34851840 & 1991911/3.014 & 34851840 & 1989953/3.017 & 34851840 \\
042 & 562868/2.997 & 8334912 & 562868/2.997 & 8334912 & 561628/3.003 & 8334912 \\
043 & 268777/2.521 & 141320 & 268777/2.521 & 141320 & 265478/2.553 & 141320 \\
044 & 96408/1.602 & 109824 & 68899/2.241 & 120384 & 67862/2.275 & 128576 \\
045 & 42601/2.949 & 0 & 42601/2.949 & 0 & 42601/2.949 & 0 \\
046 & 142176/1.944 & 4021280 & 142176/1.944 & 4021280 & 106079/2.606 & 3205536 \\
047 & 324416/1.006 & 257106 & 321355/1.016 & 1146116 & 321355/1.016 & 1146116 \\
048 & 222220/1.000 & 0 & 218425/1.017 & 592232 & 218419/1.017 & 592232 \\
049 & 160206/1.225 & 0 & 160206/1.225 & 0 & 160206/1.225 & 0 \\
050 & 108121/1.384 & 3456 & 108121/1.384 & 3456 & 107981/1.386 & 11652 \\
051 & 607628/1.000 & 0 & 607628/1.000 & 0 & 607628/1.000 & 0 \\
052 & 61071/2.952 & 181248 & 61071/2.952 & 181248 & 61071/2.952 & 181248 \\
053 & 542159/2.716 & 5803072 & 542159/2.716 & 5803072 & 491765/2.995 & 5354272 \\
054 & 517960/1.574 & 49152 & 516442/1.578 & 151622 & 513130/1.588 & 151622 \\
055 & 42047/2.826 & 115712 & 42047/2.826 & 115712 & 39987/2.972 & 123904 \\
056 & 265710/1.044 & 0 & 265710/1.044 & 0 & 265710/1.044 & 0 \\
057 & 20705/2.885 & 14112 & 20705/2.885 & 14112 & 20705/2.885 & 14112 \\
058 & 1541578/2.945 & 29859840 & 1554831/2.920 & 30400512 & 1533246/2.961 & 30445568 \\
059 & 266217/3.116 & 2998272 & 266217/3.116 & 2998272 & 257883/3.216 & 3047424 \\
060 & 298210/2.980 & 3221760 & 298210/2.980 & 3221760 & 297002/2.993 & 3225856 \\
061 & 30303/2.547 & 36096 & 30303/2.547 & 36096 & 30303/2.547 & 36096 \\
062 & 2513754/1.136 & 68020224 & 1939281/1.472 & 72442368 & 1818680/1.570 & 71553024 \\
063 & 795582/1.345 & 2116608 & 617304/1.734 & 22190592 & 593484/1.804 & 22210560 \\
064 & 19034/1.004 & 12198 & 19116/1.000 & 0 & 19116/1.000 & 0 \\
065 & 17972/2.953 & 27936 & 17972/2.953 & 27936 & 17972/2.953 & 27936 \\
066 & 248146/1.447 & 0 & 248146/1.447 & 0 & 247485/1.451 & 0 \\
067 & 23568859/2.708 & 263599744 & 23568859/2.708 & 263599744 & 22573181/2.828 & 261105536 \\
068 & 346246/1.046 & 0 & 281918/1.284 & 2187040 & 281418/1.287 & 2187040 \\
069 & 24642/1.020 & 61828 & 22905/1.098 & 0 & 22655/1.110 & 0 \\
070 & 39333/2.982 & 259072 & 39333/2.982 & 259072 & 39333/2.982 & 259072 \\
071 & 18919/1.000 & 0 & 17557/1.078 & 0 & 17401/1.087 & 0 \\
072 & 7935290/3.012 & 120023840 & 7935290/3.012 & 120023840 & 7921555/3.017 & 120056608 \\
073 & 4414910/2.284 & 139323392 & 4515227/2.233 & 140671616 & 3969025/2.540 & 139309056 \\
074 & 559584/3.006 & 7614464 & 559584/3.006 & 7614464 & 554482/3.034 & 7618560 \\
075 & 1441809/1.000 & 0 & 1363452/1.057 & 3022344 & 1363444/1.057 & 3022344 \\
076 & 2050028/3.017 & 34332672 & 2050028/3.017 & 34332672 & 2047868/3.020 & 34332672 \\
077 & 168976/2.931 & 114688 & 168976/2.931 & 114688 & 168126/2.946 & 114688 \\
078 & 993959/2.310 & 12673800 & 993959/2.310 & 12673800 & 882747/2.601 & 11678520 \\
079 & 10181692/3.050 & 39960576 & 10181692/3.050 & 39960576 & 10175315/3.052 & 39960576 \\
080 & 100325/2.912 & 3047424 & 100325/2.912 & 3047424 & 100325/2.912 & 3047424 \\
081 & 51943/2.922 & 0 & 51943/2.922 & 0 & 51943/2.922 & 0 \\
082 & 709767/1.016 & 2507492 & 619074/1.165 & 2534220 & 606773/1.188 & 2481484 \\
083 & 512623/2.064 & 14842752 & 512623/2.064 & 14842752 & 402691/2.627 & 14842752 \\
084 & 836220/2.999 & 0 & 836220/2.999 & 0 & 836220/2.999 & 0 \\
085 & 3187899/1.000 & 0 & 3130505/1.018 & 2382092 & 3112095/1.024 & 2381964 \\
086 & 118296/1.001 & 84966 & 115732/1.023 & 0 & 114476/1.034 & 0 \\
087 & 1374901/2.890 & 11036160 & 1374901/2.890 & 11036160 & 1372451/2.896 & 11036160 \\
088 & 221995/1.097 & 0 & 196677/1.239 & 2359468 & 190688/1.278 & 2359468 \\
089 & 243614/2.845 & 477440 & 243614/2.845 & 477440 & 242586/2.857 & 477440 \\
090 & 657430/1.792 & 21438504 & 657430/1.792 & 21438504 & 549431/2.145 & 5289120 \\
091 & 3078087/2.985 & 63389504 & 3078087/2.985 & 63389504 & 3057457/3.005 & 63389504 \\
092 & 2129015/2.262 & 67971680 & 2231804/2.158 & 66784736 & 1792534/2.686 & 66784736 \\
093 & 25103/2.956 & 16384 & 25103/2.956 & 16384 & 25103/2.956 & 16384 \\
094 & 32222/2.945 & 212544 & 32222/2.945 & 212544 & 31679/2.995 & 212544 \\
095 & 695775/2.996 & 0 & 695775/2.996 & 0 & 695775/2.996 & 0 \\
096 & 30318/2.777 & 79488 & 30318/2.777 & 79488 & 30318/2.777 & 79488 \\
097 & 3660459/3.139 & 16031744 & 3660459/3.139 & 16031744 & 3622881/3.172 & 16379904 \\
098 & 89756/1.977 & 0 & 89756/1.977 & 0 & 89756/1.977 & 0 \\
099 & 539533/3.008 & 8011776 & 539533/3.008 & 8011776 & 537694/3.018 & 8011776 \\
100 & 57811/2.616 & 72192 & 57811/2.616 & 72192 & 57811/2.616 & 72192 \\
\end{longtable}

\begin{longtable}{crrrrrr}
\caption{4核三问逐例工期/加速比与额外搬运（100图）}\label{tab:detail-4}\\
\toprule
图 & $T_A/S_A$ & $D_A$/B & $T_B/S_B$ & $D_B$/B & $T_C/S_C$ & $D_C$/B \\
\midrule
\endfirsthead
\multicolumn{7}{c}{续表}\\
\toprule
图 & $T_A/S_A$ & $D_A$/B & $T_B/S_B$ & $D_B$/B & $T_C/S_C$ & $D_C$/B \\
\midrule
\endhead
\midrule\multicolumn{7}{r}{续下页}\\
\endfoot
\bottomrule
\endlastfoot
001 & 58984/3.952 & 3456 & 58984/3.952 & 3456 & 58984/3.952 & 3456 \\
002 & 173068/1.514 & 594432 & 122326/2.141 & 3930624 & 119922/2.184 & 3930624 \\
003 & 538381/1.119 & 98304 & 524210/1.149 & 3647050 & 522594/1.152 & 3579466 \\
004 & 104188/3.766 & 331776 & 104188/3.766 & 331776 & 104188/3.766 & 331776 \\
005 & 95618/1.000 & 0 & 88696/1.078 & 1492040 & 86911/1.100 & 1492040 \\
006 & 31152/2.659 & 30080 & 31152/2.659 & 30080 & 31152/2.659 & 30080 \\
007 & 254398/3.980 & 82944 & 254398/3.980 & 82944 & 254398/3.980 & 82944 \\
008 & 123060/3.962 & 0 & 123060/3.962 & 0 & 123060/3.962 & 0 \\
009 & 92198/1.662 & 0 & 91278/1.679 & 9254 & 90556/1.692 & 106060 \\
010 & 29918/2.907 & 96768 & 29918/2.907 & 96768 & 29845/2.914 & 96768 \\
011 & 46880/3.940 & 540672 & 46880/3.940 & 540672 & 46880/3.940 & 540672 \\
012 & 15863/3.725 & 49536 & 15863/3.725 & 49536 & 15705/3.762 & 49536 \\
013 & 66982/3.941 & 24576 & 66982/3.941 & 24576 & 66982/3.941 & 24576 \\
014 & 4521815/3.914 & 81886104 & 4521815/3.914 & 81886104 & 4500714/3.932 & 81886104 \\
015 & 48727/3.640 & 362248 & 48727/3.640 & 362248 & 48669/3.644 & 362248 \\
016 & 7715523/1.000 & 0 & 7706495/1.001 & 4784208 & 7706472/1.001 & 4718672 \\
017 & 41209/3.951 & 253440 & 41209/3.951 & 253440 & 41209/3.951 & 253440 \\
018 & 253676/3.959 & 4938752 & 253676/3.959 & 4938752 & 253384/3.964 & 4938752 \\
019 & 24099/3.367 & 54336 & 24099/3.367 & 54336 & 23835/3.404 & 62016 \\
020 & 133704/3.985 & 0 & 133704/3.985 & 0 & 133704/3.985 & 0 \\
021 & 1550324/4.236 & 4608000 & 1550324/4.236 & 4608000 & 1550324/4.236 & 4608000 \\
022 & 14526/3.909 & 13824 & 14526/3.909 & 13824 & 14275/3.978 & 13824 \\
023 & 30142/4.068 & 281280 & 30142/4.068 & 281280 & 30142/4.068 & 281280 \\
024 & 2555343/1.000 & 0 & 2552735/1.001 & 2162724 & 2552735/1.001 & 2162724 \\
025 & 1217444/3.994 & 82944 & 1217444/3.994 & 82944 & 1217444/3.994 & 82944 \\
026 & 41665/2.381 & 95664 & 41665/2.381 & 95664 & 41665/2.381 & 95664 \\
027 & 559196/4.227 & 3317760 & 559196/4.227 & 3317760 & 559196/4.227 & 3317760 \\
028 & 13608312/3.221 & 277760640 & 13608312/3.221 & 277760640 & 13585418/3.227 & 277760640 \\
029 & 11028/3.405 & 0 & 11028/3.405 & 0 & 11028/3.405 & 0 \\
030 & 700242/4.010 & 7758848 & 700242/4.010 & 7758848 & 699694/4.013 & 7758848 \\
031 & 357987/3.951 & 7953408 & 357987/3.951 & 7953408 & 355839/3.975 & 7953408 \\
032 & 11004/3.766 & 18576 & 11004/3.766 & 18576 & 11004/3.766 & 18576 \\
033 & 443447/4.001 & 7214688 & 443447/4.001 & 7214688 & 441850/4.015 & 7219296 \\
034 & 212170/3.973 & 2978304 & 212170/3.973 & 2978304 & 212170/3.973 & 2978304 \\
035 & 95406/1.777 & 0 & 95406/1.777 & 0 & 95406/1.777 & 0 \\
036 & 233584/3.988 & 3456 & 233584/3.988 & 3456 & 233584/3.988 & 3456 \\
037 & 53320/3.939 & 614400 & 53320/3.939 & 614400 & 53320/3.939 & 614400 \\
038 & 339624/4.018 & 9295872 & 339624/4.018 & 9295872 & 337813/4.039 & 9295872 \\
039 & 395323/3.631 & 17604608 & 395323/3.631 & 17604608 & 368055/3.900 & 17604608 \\
040 & 137746/1.840 & 53760 & 137746/1.840 & 53760 & 136073/1.863 & 324618 \\
041 & 1497746/4.009 & 34397184 & 1497746/4.009 & 34397184 & 1494017/4.019 & 34397184 \\
042 & 426511/3.955 & 8553312 & 426511/3.955 & 8553312 & 423652/3.981 & 8553312 \\
043 & 268777/2.521 & 141320 & 268777/2.521 & 141320 & 265478/2.553 & 141320 \\
044 & 69401/2.225 & 1017696 & 68899/2.241 & 120384 & 67862/2.275 & 128576 \\
045 & 31878/3.940 & 0 & 31878/3.940 & 0 & 31878/3.940 & 0 \\
046 & 124612/2.219 & 4513952 & 124612/2.219 & 4513952 & 98844/2.797 & 4189760 \\
047 & 324416/1.006 & 257106 & 314757/1.037 & 2205644 & 314754/1.037 & 2205644 \\
048 & 222220/1.000 & 0 & 218425/1.017 & 592232 & 218419/1.017 & 592232 \\
049 & 160206/1.225 & 0 & 160206/1.225 & 0 & 160206/1.225 & 0 \\
050 & 108121/1.384 & 3456 & 107794/1.389 & 34380 & 107658/1.390 & 34380 \\
051 & 607628/1.000 & 0 & 541017/1.123 & 7471132 & 541017/1.123 & 7471132 \\
052 & 45610/3.953 & 254976 & 45610/3.953 & 254976 & 45610/3.953 & 254976 \\
053 & 492232/2.992 & 5772352 & 492232/2.992 & 5772352 & 491765/2.995 & 5354272 \\
054 & 517960/1.574 & 49152 & 516442/1.578 & 151622 & 513130/1.588 & 151622 \\
055 & 32099/3.702 & 173568 & 32099/3.702 & 173568 & 32099/3.702 & 173568 \\
056 & 265710/1.044 & 0 & 265710/1.044 & 0 & 265710/1.044 & 0 \\
057 & 20705/2.885 & 14112 & 16485/3.623 & 28224 & 16485/3.623 & 28224 \\
058 & 1154519/3.932 & 42622976 & 1154519/3.932 & 42622976 & 1142629/3.973 & 42622976 \\
059 & 185216/4.478 & 1081344 & 185216/4.478 & 1081344 & 185216/4.478 & 1081344 \\
060 & 221720/4.009 & 3930376 & 221720/4.009 & 3930376 & 221694/4.009 & 3656968 \\
061 & 24024/3.213 & 70656 & 24024/3.213 & 70656 & 24024/3.213 & 70656 \\
062 & 2377293/1.201 & 70643712 & 1767455/1.615 & 67070976 & 1270730/2.247 & 6379008 \\
063 & 348902/3.068 & 2469888 & 324251/3.301 & 2211840 & 324251/3.301 & 2211840 \\
064 & 19034/1.004 & 12198 & 19116/1.000 & 0 & 19116/1.000 & 0 \\
065 & 13244/4.007 & 41760 & 13244/4.007 & 41760 & 13062/4.063 & 41760 \\
066 & 248146/1.447 & 0 & 248146/1.447 & 0 & 247485/1.451 & 0 \\
067 & 18473155/3.456 & 263174912 & 18473155/3.456 & 263174912 & 17693167/3.608 & 261383424 \\
068 & 346246/1.046 & 0 & 267505/1.354 & 2421242 & 266974/1.356 & 2421242 \\
069 & 24642/1.020 & 61828 & 22905/1.098 & 0 & 22655/1.110 & 0 \\
070 & 31225/3.757 & 360960 & 31225/3.757 & 360960 & 31192/3.761 & 360960 \\
071 & 18919/1.000 & 0 & 17557/1.078 & 0 & 17401/1.087 & 0 \\
072 & 5993260/3.987 & 122669856 & 5993260/3.987 & 122669856 & 5988828/3.990 & 122669856 \\
073 & 3903787/2.583 & 140222208 & 3903787/2.583 & 140222208 & 3434743/2.935 & 139257088 \\
074 & 418987/4.015 & 7852032 & 418987/4.015 & 7852032 & 415633/4.047 & 7852032 \\
075 & 1441809/1.000 & 0 & 1320997/1.091 & 4283756 & 1320971/1.091 & 4283756 \\
076 & 1540201/4.016 & 37442560 & 1540201/4.016 & 37442560 & 1536296/4.026 & 37453824 \\
077 & 127832/3.875 & 114688 & 127832/3.875 & 114688 & 126652/3.911 & 131072 \\
078 & 619540/3.706 & 15039360 & 690924/3.323 & 13146912 & 642020/3.576 & 12673464 \\
079 & 7635624/4.067 & 30445568 & 7635624/4.067 & 30445568 & 7635624/4.067 & 30445568 \\
080 & 100325/2.912 & 3047424 & 100325/2.912 & 3047424 & 90559/3.226 & 4190208 \\
081 & 39588/3.834 & 0 & 39588/3.834 & 0 & 39588/3.834 & 0 \\
082 & 709767/1.016 & 2507492 & 581815/1.239 & 5394700 & 581598/1.240 & 5394700 \\
083 & 450029/2.351 & 14977600 & 450029/2.351 & 14977600 & 333151/3.176 & 14977600 \\
084 & 629028/3.986 & 0 & 629028/3.986 & 0 & 629028/3.986 & 0 \\
085 & 3187899/1.000 & 0 & 2987335/1.067 & 11728024 & 2986917/1.067 & 11728024 \\
086 & 109864/1.078 & 439126 & 94047/1.259 & 398416 & 88995/1.330 & 398416 \\
087 & 1046846/3.796 & 12607488 & 1046846/3.796 & 12607488 & 1042446/3.812 & 12612096 \\
088 & 221995/1.097 & 0 & 185641/1.312 & 2459964 & 184608/1.320 & 2490684 \\
089 & 178160/3.890 & 716160 & 178160/3.890 & 716160 & 175147/3.957 & 716160 \\
090 & 608275/1.937 & 21392136 & 608275/1.937 & 21392136 & 534087/2.206 & 20329752 \\
091 & 2314405/3.969 & 66813600 & 2314405/3.969 & 66813600 & 2306119/3.984 & 66813600 \\
092 & 1927728/2.498 & 67180384 & 1927728/2.498 & 67180384 & 1470508/3.275 & 67180384 \\
093 & 19779/3.752 & 24576 & 19779/3.752 & 24576 & 19779/3.752 & 24576 \\
094 & 24640/3.851 & 318528 & 24640/3.851 & 318528 & 24640/3.851 & 318528 \\
095 & 524422/3.976 & 0 & 524422/3.976 & 0 & 524422/3.976 & 0 \\
096 & 26984/3.120 & 79488 & 26984/3.120 & 79488 & 26984/3.120 & 79488 \\
097 & 2626908/4.375 & 14221312 & 2626908/4.375 & 14221312 & 2626908/4.375 & 14221312 \\
098 & 89756/1.977 & 0 & 89756/1.977 & 0 & 89756/1.977 & 0 \\
099 & 405336/4.004 & 9838592 & 405336/4.004 & 9838592 & 403147/4.026 & 9256960 \\
100 & 57811/2.616 & 72192 & 57811/2.616 & 72192 & 57811/2.616 & 72192 \\
\end{longtable}

\begin{longtable}{crrrrrr}
\caption{5核三问逐例工期/加速比与额外搬运（100图）}\label{tab:detail-5}\\
\toprule
图 & $T_A/S_A$ & $D_A$/B & $T_B/S_B$ & $D_B$/B & $T_C/S_C$ & $D_C$/B \\
\midrule
\endfirsthead
\multicolumn{7}{c}{续表}\\
\toprule
图 & $T_A/S_A$ & $D_A$/B & $T_B/S_B$ & $D_B$/B & $T_C/S_C$ & $D_C$/B \\
\midrule
\endhead
\midrule\multicolumn{7}{r}{续下页}\\
\endfoot
\bottomrule
\endlastfoot
001 & 47502/4.907 & 4608 & 47502/4.907 & 4608 & 47502/4.907 & 4608 \\
002 & 128412/2.040 & 700416 & 111678/2.346 & 4024320 & 105055/2.493 & 4010496 \\
003 & 538381/1.119 & 98304 & 469702/1.282 & 5272020 & 469516/1.283 & 5272020 \\
004 & 88341/4.441 & 442368 & 88341/4.441 & 442368 & 86005/4.562 & 473088 \\
005 & 95618/1.000 & 0 & 88696/1.078 & 1492040 & 86911/1.100 & 1492040 \\
006 & 31152/2.659 & 30080 & 31152/2.659 & 30080 & 31152/2.659 & 30080 \\
007 & 203834/4.968 & 110592 & 203834/4.968 & 110592 & 203834/4.968 & 110592 \\
008 & 100603/4.847 & 0 & 100603/4.847 & 0 & 100603/4.847 & 0 \\
009 & 92198/1.662 & 0 & 91278/1.679 & 9254 & 90556/1.692 & 106060 \\
010 & 29918/2.907 & 96768 & 29918/2.907 & 96768 & 29845/2.914 & 96768 \\
011 & 46880/3.940 & 540672 & 46880/3.940 & 540672 & 46777/3.948 & 1622016 \\
012 & 13406/4.408 & 53568 & 13406/4.408 & 53568 & 13406/4.408 & 53568 \\
013 & 54363/4.855 & 32768 & 54349/4.857 & 49152 & 53418/4.941 & 49152 \\
014 & 3681597/4.807 & 87153672 & 3681597/4.807 & 87153672 & 3672221/4.820 & 87153672 \\
015 & 43665/4.062 & 336384 & 43665/4.062 & 336384 & 42872/4.137 & 336384 \\
016 & 7715523/1.000 & 0 & 7706495/1.001 & 4784208 & 7706472/1.001 & 4718672 \\
017 & 34231/4.756 & 338688 & 34231/4.756 & 338688 & 34231/4.756 & 338688 \\
018 & 205219/4.894 & 5345280 & 205219/4.894 & 5345280 & 201659/4.981 & 5345280 \\
019 & 24099/3.367 & 54336 & 24099/3.367 & 54336 & 23835/3.404 & 62016 \\
020 & 107650/4.949 & 0 & 107650/4.949 & 0 & 107650/4.949 & 0 \\
021 & 1243902/5.280 & 7372800 & 1250168/5.254 & 11431936 & 1249380/5.257 & 11431936 \\
022 & 11824/4.802 & 20736 & 11824/4.802 & 20736 & 11824/4.802 & 20736 \\
023 & 24517/5.001 & 270912 & 24517/5.001 & 270912 & 24517/5.001 & 270912 \\
024 & 2555343/1.000 & 0 & 2538844/1.006 & 6291556 & 2538844/1.006 & 6291556 \\
025 & 975282/4.986 & 110592 & 974747/4.989 & 118272 & 974747/4.989 & 118272 \\
026 & 41665/2.381 & 95664 & 41665/2.381 & 95664 & 41665/2.381 & 95664 \\
027 & 558753/4.231 & 5529600 & 557147/4.243 & 5533696 & 552561/4.278 & 5533696 \\
028 & 11491035/3.815 & 277818912 & 11491035/3.815 & 277818912 & 10636977/4.121 & 262763520 \\
029 & 9383/4.002 & 0 & 9383/4.002 & 0 & 9383/4.002 & 0 \\
030 & 559889/5.015 & 10750976 & 559889/5.015 & 10750976 & 559025/5.023 & 10760192 \\
031 & 285693/4.951 & 7931904 & 285693/4.951 & 7931904 & 284489/4.972 & 7931904 \\
032 & 10496/3.948 & 12432 & 10496/3.948 & 12432 & 10479/3.955 & 12432 \\
033 & 355102/4.996 & 7816320 & 355102/4.996 & 7816320 & 354514/5.005 & 7816320 \\
034 & 170371/4.948 & 3621888 & 170371/4.948 & 3621888 & 168887/4.991 & 3631104 \\
035 & 95406/1.777 & 0 & 95406/1.777 & 0 & 95406/1.777 & 0 \\
036 & 187182/4.976 & 4608 & 187182/4.976 & 4608 & 187182/4.976 & 4608 \\
037 & 43800/4.795 & 1925120 & 43800/4.795 & 1925120 & 43013/4.882 & 1925120 \\
038 & 275133/4.960 & 9676800 & 275133/4.960 & 9676800 & 272533/5.007 & 9676800 \\
039 & 316542/4.534 & 9805824 & 395323/3.631 & 17604608 & 340476/4.215 & 20480000 \\
040 & 137746/1.840 & 53760 & 137746/1.840 & 53760 & 136073/1.863 & 324618 \\
041 & 1201884/4.995 & 40673280 & 1201884/4.995 & 40673280 & 1194643/5.026 & 40673280 \\
042 & 342520/4.925 & 10187328 & 342520/4.925 & 10187328 & 339457/4.969 & 10187328 \\
043 & 268777/2.521 & 141320 & 268777/2.521 & 141320 & 265478/2.553 & 141320 \\
044 & 69401/2.225 & 1017696 & 68899/2.241 & 120384 & 63439/2.434 & 143616 \\
045 & 26164/4.801 & 0 & 26164/4.801 & 0 & 26164/4.801 & 0 \\
046 & 108698/2.543 & 4646464 & 117204/2.359 & 4158144 & 98844/2.797 & 4189760 \\
047 & 324416/1.006 & 257106 & 314757/1.037 & 2205644 & 314754/1.037 & 2205644 \\
048 & 222220/1.000 & 0 & 217218/1.023 & 717070 & 215393/1.032 & 717070 \\
049 & 160206/1.225 & 0 & 160206/1.225 & 0 & 160206/1.225 & 0 \\
050 & 108121/1.384 & 3456 & 107794/1.389 & 34380 & 107658/1.390 & 34380 \\
051 & 607628/1.000 & 0 & 541017/1.123 & 7471132 & 541017/1.123 & 7471132 \\
052 & 37387/4.823 & 328704 & 37387/4.823 & 328704 & 37387/4.823 & 328704 \\
053 & 491765/2.995 & 4469536 & 492232/2.992 & 5772352 & 491765/2.995 & 5354272 \\
054 & 517960/1.574 & 49152 & 516442/1.578 & 151622 & 513130/1.588 & 151622 \\
055 & 25635/4.636 & 231424 & 25635/4.636 & 231424 & 25635/4.636 & 231424 \\
056 & 265710/1.044 & 0 & 265710/1.044 & 0 & 265710/1.044 & 0 \\
057 & 14028/4.258 & 14112 & 14028/4.258 & 14112 & 14028/4.258 & 14112 \\
058 & 1004819/4.518 & 32104448 & 1004819/4.518 & 32104448 & 1004819/4.518 & 32104448 \\
059 & 150782/5.501 & 4325376 & 150782/5.501 & 4325376 & 150736/5.503 & 4325376 \\
060 & 175702/5.059 & 3948928 & 175702/5.059 & 3948928 & 175626/5.061 & 3948928 \\
061 & 22283/3.464 & 82944 & 22283/3.464 & 82944 & 22283/3.464 & 82944 \\
062 & 1661193/1.719 & 7405056 & 719016/3.970 & 7147008 & 701554/4.069 & 9028608 \\
063 & 269453/3.972 & 2755584 & 251581/4.255 & 2509824 & 251581/4.255 & 2509824 \\
064 & 19034/1.004 & 12198 & 19116/1.000 & 0 & 19116/1.000 & 0 \\
065 & 11396/4.657 & 55872 & 11396/4.657 & 55872 & 11396/4.657 & 55872 \\
066 & 248146/1.447 & 0 & 237259/1.513 & 3464460 & 237121/1.514 & 3464460 \\
067 & 15931687/4.007 & 262750080 & 15931687/4.007 & 262750080 & 15229262/4.192 & 261427072 \\
068 & 346246/1.046 & 0 & 267505/1.354 & 2421242 & 266974/1.356 & 2421242 \\
069 & 22097/1.138 & 259792 & 13375/1.880 & 271980 & 13239/1.899 & 271980 \\
070 & 26275/4.464 & 449024 & 26275/4.464 & 449024 & 26030/4.506 & 455936 \\
071 & 18919/1.000 & 0 & 17557/1.078 & 0 & 17401/1.087 & 0 \\
072 & 4831281/4.946 & 130884272 & 4831281/4.946 & 130884272 & 4809827/4.969 & 130884272 \\
073 & 3517515/2.866 & 139772800 & 3517515/2.866 & 139772800 & 3434743/2.935 & 139257088 \\
074 & 337275/4.988 & 7983104 & 337275/4.988 & 7983104 & 335626/5.012 & 7983104 \\
075 & 1441809/1.000 & 0 & 1219033/1.183 & 4890512 & 1219026/1.183 & 4890512 \\
076 & 1239374/4.991 & 38022144 & 1239374/4.991 & 38022144 & 1232497/5.018 & 38022144 \\
077 & 115617/4.284 & 172032 & 115617/4.284 & 172032 & 115617/4.284 & 172032 \\
078 & 619540/3.706 & 15039360 & 663941/3.458 & 13620024 & 578013/3.972 & 13620024 \\
079 & 6075976/5.111 & 47677440 & 6075976/5.111 & 47677440 & 6075472/5.111 & 47677440 \\
080 & 100325/2.912 & 3047424 & 100325/2.912 & 3047424 & 88140/3.314 & 4603904 \\
081 & 30943/4.906 & 0 & 30943/4.906 & 0 & 30943/4.906 & 0 \\
082 & 709767/1.016 & 2507492 & 581815/1.239 & 5394700 & 581598/1.240 & 5394700 \\
083 & 442480/2.391 & 15058048 & 442480/2.391 & 15058048 & 310731/3.405 & 15058048 \\
084 & 503472/4.980 & 0 & 503472/4.980 & 0 & 503472/4.980 & 0 \\
085 & 3187899/1.000 & 0 & 2791092/1.142 & 13660374 & 2986917/1.067 & 11728024 \\
086 & 109864/1.078 & 439126 & 94047/1.259 & 398416 & 88995/1.330 & 398416 \\
087 & 833349/4.769 & 12487680 & 833349/4.769 & 12487680 & 829449/4.791 & 12487680 \\
088 & 221995/1.097 & 0 & 185641/1.312 & 2459964 & 184608/1.320 & 2490684 \\
089 & 147432/4.701 & 954880 & 147432/4.701 & 954880 & 144388/4.800 & 954880 \\
090 & 542638/2.172 & 21724704 & 542638/2.172 & 21724704 & 416535/2.829 & 21213240 \\
091 & 1882479/4.880 & 73417216 & 1882479/4.880 & 73417216 & 1876653/4.895 & 73417216 \\
092 & 1661234/2.899 & 66980256 & 1649430/2.920 & 66980256 & 1351088/3.564 & 66980256 \\
093 & 15876/4.674 & 32768 & 15876/4.674 & 32768 & 15876/4.674 & 32768 \\
094 & 21573/4.399 & 424800 & 21573/4.399 & 424800 & 20380/4.656 & 424800 \\
095 & 420852/4.954 & 0 & 420852/4.954 & 0 & 420852/4.954 & 0 \\
096 & 26984/3.120 & 79488 & 26984/3.120 & 79488 & 26984/3.120 & 79488 \\
097 & 2102348/5.466 & 21311488 & 2102348/5.466 & 21311488 & 2102250/5.466 & 21311488 \\
098 & 89756/1.977 & 0 & 89756/1.977 & 0 & 89756/1.977 & 0 \\
099 & 336826/4.819 & 10371072 & 336826/4.819 & 10371072 & 327139/4.961 & 10343424 \\
100 & 57811/2.616 & 72192 & 57811/2.616 & 72192 & 57811/2.616 & 72192 \\
\end{longtable}

\endgroup

\section{Cache同计划完整配对}\label{app:cache}
表\ref{tab:cache-detail-1}至表\ref{tab:cache-detail-5}共列出500组配对。
$T_B=F_B(X_B^{(K)})$，$T_C(X_B)=F_C(X_B^{(K)})$，
$T_C(X_C)=F_C(X_C^{(K)})$；三比值由式\eqref{eq:cache-factors}计算。
表中同时列出$C(X_B)$与$C(X_C)$各自的字节命中率；硬件效应图和同计划工期解释只使用$C(X_B)$字段。
每组的$C(X_B)$沿用该组$B$的最终计划，因而三段工期的计划身份保持一致。
\begingroup
\scriptsize
\setlength{\tabcolsep}{2pt}
\setlength{\LTleft}{\fill}\setlength{\LTright}{\fill}
\renewcommand{\arraystretch}{1.03}
\begin{longtable}{crrrrrrrr}
\caption{1核Cache配对（100图）}\label{tab:cache-detail-1}\\
\toprule
图 & $T_B$ & $T_C(X_B)$ & $T_C(X_C)$ & $S_{hw}$ & $S_{ad}$ & $S_{opt}$ & $C(X_B)$命中/\% & $C(X_C)$命中/\% \\
\midrule
\endfirsthead
\multicolumn{9}{c}{续表}\\
\toprule
图 & $T_B$ & $T_C(X_B)$ & $T_C(X_C)$ & $S_{hw}$ & $S_{ad}$ & $S_{opt}$ & $C(X_B)$命中/\% & $C(X_C)$命中/\% \\
\midrule
\endhead
\midrule\multicolumn{9}{r}{续下页}\\
\endfoot
\bottomrule
\endlastfoot
001 & 233110 & 233110 & 233110 & 1.0000 & 1.0000 & 1.0000 & 0.00 & 0.00 \\
002 & 260509 & 254213 & 254213 & 1.0248 & 1.0000 & 1.0248 & 37.89 & 37.89 \\
003 & 592981 & 592981 & 591743 & 1.0000 & 1.0021 & 1.0021 & 0.00 & 0.00 \\
004 & 392195 & 392195 & 387184 & 1.0000 & 1.0129 & 1.0129 & 22.04 & 22.69 \\
005 & 93310 & 93310 & 93171 & 1.0000 & 1.0015 & 1.0015 & 0.00 & 0.00 \\
006 & 73949 & 73949 & 73949 & 1.0000 & 1.0000 & 1.0000 & 1.19 & 1.19 \\
007 & 1012562 & 1012562 & 1012562 & 1.0000 & 1.0000 & 1.0000 & 0.00 & 0.00 \\
008 & 269325 & 269325 & 269325 & 1.0000 & 1.0000 & 1.0000 & 0.00 & 0.00 \\
009 & 153201 & 153201 & 153201 & 1.0000 & 1.0000 & 1.0000 & 0.00 & 0.00 \\
010 & 81326 & 81326 & 79774 & 1.0000 & 1.0195 & 1.0195 & 0.00 & 0.00 \\
011 & 184697 & 184697 & 184697 & 1.0000 & 1.0000 & 1.0000 & 0.00 & 0.00 \\
012 & 42251 & 42251 & 42239 & 1.0000 & 1.0003 & 1.0003 & 0.00 & 0.00 \\
013 & 263949 & 263949 & 263949 & 1.0000 & 1.0000 & 1.0000 & 0.00 & 0.00 \\
014 & 17698626 & 17689602 & 17689602 & 1.0005 & 1.0000 & 1.0005 & 12.29 & 12.29 \\
015 & 167767 & 167630 & 167630 & 1.0008 & 1.0000 & 1.0008 & 16.50 & 16.50 \\
016 & 7715523 & 7715523 & 7715523 & 1.0000 & 1.0000 & 1.0000 & 0.00 & 0.00 \\
017 & 149759 & 149759 & 149367 & 1.0000 & 1.0026 & 1.0026 & 3.54 & 3.54 \\
018 & 990647 & 990647 & 990647 & 1.0000 & 1.0000 & 1.0000 & 1.67 & 1.67 \\
019 & 70264 & 70264 & 70264 & 1.0000 & 1.0000 & 1.0000 & 0.00 & 0.00 \\
020 & 532746 & 532746 & 532746 & 1.0000 & 1.0000 & 1.0000 & 0.00 & 0.00 \\
021 & 6567937 & 6558837 & 6558837 & 1.0014 & 1.0000 & 1.0014 & 50.06 & 50.06 \\
022 & 54723 & 54723 & 54591 & 1.0000 & 1.0024 & 1.0024 & 6.54 & 11.09 \\
023 & 99763 & 99763 & 93197 & 1.0000 & 1.0705 & 1.0705 & 1.29 & 1.29 \\
024 & 2555343 & 2555343 & 2555343 & 1.0000 & 1.0000 & 1.0000 & 0.00 & 0.00 \\
025 & 4863026 & 4863026 & 4863026 & 1.0000 & 1.0000 & 1.0000 & 0.00 & 0.00 \\
026 & 95759 & 95048 & 88492 & 1.0075 & 1.0741 & 1.0821 & 13.53 & 10.63 \\
027 & 2304836 & 2304836 & 2304836 & 1.0000 & 1.0000 & 1.0000 & 5.89 & 5.89 \\
028 & 42901220 & 41700132 & 41700132 & 1.0288 & 1.0000 & 1.0288 & 63.66 & 63.66 \\
029 & 37555 & 37555 & 37555 & 1.0000 & 1.0000 & 1.0000 & 0.00 & 0.00 \\
030 & 2808034 & 2807891 & 2806776 & 1.0001 & 1.0004 & 1.0004 & 25.99 & 28.61 \\
031 & 1414369 & 1414159 & 1414159 & 1.0001 & 1.0000 & 1.0001 & 26.06 & 26.06 \\
032 & 34141 & 34141 & 34141 & 1.0000 & 1.0000 & 1.0000 & 3.54 & 0.00 \\
033 & 1774196 & 1774080 & 1774080 & 1.0001 & 1.0000 & 1.0001 & 9.63 & 9.63 \\
034 & 842996 & 842733 & 841851 & 1.0003 & 1.0010 & 1.0014 & 30.90 & 20.29 \\
035 & 169168 & 169168 & 168511 & 1.0000 & 1.0039 & 1.0039 & 0.00 & 0.00 \\
036 & 931510 & 931510 & 931510 & 1.0000 & 1.0000 & 1.0000 & 0.00 & 0.00 \\
037 & 210001 & 210001 & 210001 & 1.0000 & 1.0000 & 1.0000 & 0.00 & 0.00 \\
038 & 1364555 & 1363312 & 1359544 & 1.0009 & 1.0028 & 1.0037 & 17.36 & 28.97 \\
039 & 1435262 & 1435262 & 1426335 & 1.0000 & 1.0063 & 1.0063 & 0.00 & 0.33 \\
040 & 252193 & 252089 & 249653 & 1.0004 & 1.0098 & 1.0102 & 0.83 & 0.83 \\
041 & 6003957 & 6001575 & 6001575 & 1.0004 & 1.0000 & 1.0004 & 10.16 & 10.16 \\
042 & 1686740 & 1675633 & 1675633 & 1.0066 & 1.0000 & 1.0066 & 31.87 & 31.87 \\
043 & 657411 & 657372 & 656100 & 1.0001 & 1.0019 & 1.0020 & 4.99 & 4.99 \\
044 & 88538 & 88538 & 88538 & 1.0000 & 1.0000 & 1.0000 & 0.00 & 0.00 \\
045 & 125614 & 125614 & 125614 & 1.0000 & 1.0000 & 1.0000 & 0.00 & 0.00 \\
046 & 226233 & 226233 & 226233 & 1.0000 & 1.0000 & 1.0000 & 0.00 & 0.00 \\
047 & 323597 & 323597 & 322738 & 1.0000 & 1.0027 & 1.0027 & 0.00 & 0.00 \\
048 & 222220 & 222220 & 222220 & 1.0000 & 1.0000 & 1.0000 & 0.00 & 0.00 \\
049 & 195430 & 195430 & 194548 & 1.0000 & 1.0045 & 1.0045 & 0.00 & 0.00 \\
050 & 149242 & 149242 & 148876 & 1.0000 & 1.0025 & 1.0025 & 4.58 & 3.28 \\
051 & 607628 & 607628 & 607628 & 1.0000 & 1.0000 & 1.0000 & 0.00 & 0.00 \\
052 & 159895 & 159895 & 158577 & 1.0000 & 1.0083 & 1.0083 & 0.00 & 0.00 \\
053 & 1448004 & 1426552 & 1421619 & 1.0150 & 1.0035 & 1.0186 & 39.64 & 39.00 \\
054 & 796718 & 796718 & 792519 & 1.0000 & 1.0053 & 1.0053 & 0.00 & 0.00 \\
055 & 118835 & 118835 & 114628 & 1.0000 & 1.0367 & 1.0367 & 10.01 & 4.89 \\
056 & 277285 & 277285 & 277285 & 1.0000 & 1.0000 & 1.0000 & 0.00 & 0.00 \\
057 & 59247 & 59247 & 59096 & 1.0000 & 1.0026 & 1.0026 & 13.14 & 9.08 \\
058 & 4539717 & 4539717 & 4533133 & 1.0000 & 1.0015 & 1.0015 & 0.00 & 0.00 \\
059 & 829473 & 815938 & 815938 & 1.0166 & 1.0000 & 1.0166 & 61.00 & 61.00 \\
060 & 885660 & 884971 & 883441 & 1.0008 & 1.0017 & 1.0025 & 25.49 & 24.15 \\
061 & 77191 & 77191 & 77191 & 1.0000 & 1.0000 & 1.0000 & 0.00 & 0.00 \\
062 & 2854809 & 2854809 & 2854809 & 1.0000 & 1.0000 & 1.0000 & 0.00 & 0.00 \\
063 & 1070397 & 1070397 & 1070397 & 1.0000 & 1.0000 & 1.0000 & 0.00 & 0.00 \\
064 & 19116 & 19116 & 19116 & 1.0000 & 1.0000 & 1.0000 & 0.00 & 0.00 \\
065 & 39901 & 39901 & 39824 & 1.0000 & 1.0019 & 1.0019 & 0.00 & 0.00 \\
066 & 357763 & 357763 & 356326 & 1.0000 & 1.0040 & 1.0040 & 0.00 & 0.00 \\
067 & 63397576 & 63397576 & 63397576 & 1.0000 & 1.0000 & 1.0000 & 0.00 & 0.00 \\
068 & 360298 & 360298 & 359166 & 1.0000 & 1.0032 & 1.0032 & 0.00 & 0.00 \\
069 & 22905 & 22905 & 22655 & 1.0000 & 1.0110 & 1.0110 & 0.00 & 0.00 \\
070 & 115088 & 115088 & 104334 & 1.0000 & 1.1031 & 1.1031 & 0.00 & 0.00 \\
071 & 17557 & 17557 & 17401 & 1.0000 & 1.0090 & 1.0090 & 0.00 & 0.00 \\
072 & 23897828 & 23890756 & 23890756 & 1.0003 & 1.0000 & 1.0003 & 2.62 & 2.62 \\
073 & 9870212 & 9870212 & 9870212 & 1.0000 & 1.0000 & 1.0000 & 0.00 & 0.00 \\
074 & 1682239 & 1678120 & 1678120 & 1.0025 & 1.0000 & 1.0025 & 16.68 & 16.68 \\
075 & 1433562 & 1433562 & 1427745 & 1.0000 & 1.0041 & 1.0041 & 0.00 & 0.00 \\
076 & 6185113 & 6183098 & 6182761 & 1.0003 & 1.0001 & 1.0004 & 13.36 & 14.12 \\
077 & 495320 & 495320 & 494044 & 1.0000 & 1.0026 & 1.0026 & 0.00 & 0.00 \\
078 & 2232856 & 2188693 & 2142955 & 1.0202 & 1.0213 & 1.0420 & 55.23 & 59.38 \\
079 & 31052251 & 31052251 & 31052251 & 1.0000 & 1.0000 & 1.0000 & 0.00 & 0.00 \\
080 & 292138 & 291826 & 291826 & 1.0011 & 1.0000 & 1.0011 & 16.86 & 16.86 \\
081 & 151795 & 151795 & 151795 & 1.0000 & 1.0000 & 1.0000 & 0.00 & 0.00 \\
082 & 715588 & 715588 & 709534 & 1.0000 & 1.0085 & 1.0085 & 0.00 & 0.00 \\
083 & 876669 & 876669 & 875735 & 1.0000 & 1.0011 & 1.0011 & 0.00 & 0.00 \\
084 & 1371156 & 1371156 & 1371156 & 1.0000 & 1.0000 & 1.0000 & 0.00 & 0.00 \\
085 & 3177963 & 3177963 & 3173461 & 1.0000 & 1.0014 & 1.0014 & 0.00 & 0.00 \\
086 & 115732 & 115732 & 114476 & 1.0000 & 1.0110 & 1.0110 & 0.00 & 0.00 \\
087 & 3974071 & 3969779 & 3969779 & 1.0011 & 1.0000 & 1.0011 & 18.34 & 18.34 \\
088 & 243542 & 243542 & 242071 & 1.0000 & 1.0061 & 1.0061 & 0.00 & 0.00 \\
089 & 689572 & 689572 & 689572 & 1.0000 & 1.0000 & 1.0000 & 0.00 & 0.00 \\
090 & 1087442 & 989385 & 919256 & 1.0991 & 1.0763 & 1.1830 & 61.19 & 73.32 \\
091 & 9186769 & 9168820 & 9166672 & 1.0020 & 1.0002 & 1.0022 & 10.54 & 14.99 \\
092 & 4010291 & 4010291 & 4004200 & 1.0000 & 1.0015 & 1.0015 & 0.00 & 0.00 \\
093 & 74205 & 74205 & 74205 & 1.0000 & 1.0000 & 1.0000 & 0.00 & 0.00 \\
094 & 91195 & 91195 & 87488 & 1.0000 & 1.0424 & 1.0424 & 3.03 & 3.03 \\
095 & 2070244 & 2070244 & 1161033 & 1.0000 & 1.7831 & 1.7831 & 0.00 & 0.00 \\
096 & 77108 & 77108 & 75108 & 1.0000 & 1.0266 & 1.0266 & 0.00 & 0.00 \\
097 & 11491509 & 11491509 & 11491509 & 1.0000 & 1.0000 & 1.0000 & 0.00 & 0.00 \\
098 & 176752 & 176752 & 174782 & 1.0000 & 1.0113 & 1.0113 & 2.30 & 2.72 \\
099 & 1622999 & 1618902 & 1618902 & 1.0025 & 1.0000 & 1.0025 & 25.24 & 25.24 \\
100 & 148885 & 148885 & 142753 & 1.0000 & 1.0430 & 1.0430 & 0.00 & 0.00 \\
\end{longtable}

\begin{longtable}{crrrrrrrr}
\caption{2核Cache配对（100图）}\label{tab:cache-detail-2}\\
\toprule
图 & $T_B$ & $T_C(X_B)$ & $T_C(X_C)$ & $S_{hw}$ & $S_{ad}$ & $S_{opt}$ & $C(X_B)$命中/\% & $C(X_C)$命中/\% \\
\midrule
\endfirsthead
\multicolumn{9}{c}{续表}\\
\toprule
图 & $T_B$ & $T_C(X_B)$ & $T_C(X_C)$ & $S_{hw}$ & $S_{ad}$ & $S_{opt}$ & $C(X_B)$命中/\% & $C(X_C)$命中/\% \\
\midrule
\endhead
\midrule\multicolumn{9}{r}{续下页}\\
\endfoot
\bottomrule
\endlastfoot
001 & 116868 & 116868 & 116868 & 1.0000 & 1.0000 & 1.0000 & 0.00 & 0.00 \\
002 & 165091 & 164754 & 164754 & 1.0020 & 1.0000 & 1.0020 & 18.31 & 18.31 \\
003 & 538381 & 538381 & 531381 & 1.0000 & 1.0132 & 1.0132 & 4.74 & 7.03 \\
004 & 199074 & 199074 & 199074 & 1.0000 & 1.0000 & 1.0000 & 0.00 & 0.00 \\
005 & 93310 & 93310 & 93060 & 1.0000 & 1.0027 & 1.0027 & 0.00 & 5.54 \\
006 & 42696 & 42696 & 41180 & 1.0000 & 1.0368 & 1.0368 & 9.37 & 9.37 \\
007 & 507494 & 507494 & 507494 & 1.0000 & 1.0000 & 1.0000 & 0.00 & 0.00 \\
008 & 244266 & 244266 & 244266 & 1.0000 & 1.0000 & 1.0000 & 0.00 & 0.00 \\
009 & 92198 & 92198 & 92198 & 1.0000 & 1.0000 & 1.0000 & 0.00 & 0.00 \\
010 & 45272 & 45272 & 41774 & 1.0000 & 1.0837 & 1.0837 & 11.19 & 11.19 \\
011 & 92646 & 92646 & 92646 & 1.0000 & 1.0000 & 1.0000 & 0.00 & 0.00 \\
012 & 30058 & 30058 & 29156 & 1.0000 & 1.0309 & 1.0309 & 5.72 & 5.72 \\
013 & 132198 & 132198 & 132198 & 1.0000 & 1.0000 & 1.0000 & 0.00 & 0.00 \\
014 & 8897730 & 8888616 & 8888616 & 1.0010 & 1.0000 & 1.0010 & 22.81 & 22.81 \\
015 & 90509 & 90509 & 87390 & 1.0000 & 1.0357 & 1.0357 & 18.44 & 18.44 \\
016 & 7712915 & 7712915 & 7712915 & 1.0000 & 1.0000 & 1.0000 & 0.00 & 0.00 \\
017 & 81772 & 81772 & 81772 & 1.0000 & 1.0000 & 1.0000 & 1.30 & 1.30 \\
018 & 504468 & 502571 & 502571 & 1.0038 & 1.0000 & 1.0038 & 50.57 & 50.57 \\
019 & 41382 & 41382 & 39509 & 1.0000 & 1.0474 & 1.0474 & 13.16 & 13.16 \\
020 & 266580 & 266580 & 266580 & 1.0000 & 1.0000 & 1.0000 & 0.00 & 0.00 \\
021 & 3194570 & 3193634 & 3193634 & 1.0003 & 1.0000 & 1.0003 & 21.78 & 21.78 \\
022 & 28324 & 28324 & 28324 & 1.0000 & 1.0000 & 1.0000 & 3.14 & 3.14 \\
023 & 61177 & 61177 & 60171 & 1.0000 & 1.0167 & 1.0167 & 16.23 & 16.23 \\
024 & 2552735 & 2552735 & 2552735 & 1.0000 & 1.0000 & 1.0000 & 0.00 & 0.00 \\
025 & 2431840 & 2431840 & 2431840 & 1.0000 & 1.0000 & 1.0000 & 0.00 & 0.00 \\
026 & 55057 & 55057 & 51429 & 1.0000 & 1.0705 & 1.0705 & 4.29 & 4.29 \\
027 & 1120962 & 1120962 & 1120962 & 1.0000 & 1.0000 & 1.0000 & 8.90 & 8.90 \\
028 & 23638628 & 23615145 & 23615145 & 1.0010 & 1.0000 & 1.0010 & 0.47 & 0.47 \\
029 & 19740 & 19740 & 19740 & 1.0000 & 1.0000 & 1.0000 & 0.00 & 0.00 \\
030 & 1399149 & 1398912 & 1398912 & 1.0002 & 1.0000 & 1.0002 & 41.52 & 41.52 \\
031 & 712596 & 711944 & 706497 & 1.0009 & 1.0077 & 1.0086 & 43.05 & 25.78 \\
032 & 21040 & 21040 & 20711 & 1.0000 & 1.0159 & 1.0159 & 2.18 & 2.18 \\
033 & 889697 & 887951 & 887951 & 1.0020 & 1.0000 & 1.0020 & 35.16 & 35.16 \\
034 & 418129 & 417777 & 417777 & 1.0008 & 1.0000 & 1.0008 & 41.14 & 41.14 \\
035 & 95406 & 95406 & 95406 & 1.0000 & 1.0000 & 1.0000 & 0.00 & 0.00 \\
036 & 466068 & 466068 & 466068 & 1.0000 & 1.0000 & 1.0000 & 0.00 & 0.00 \\
037 & 105374 & 105374 & 105374 & 1.0000 & 1.0000 & 1.0000 & 0.00 & 0.00 \\
038 & 674500 & 674382 & 674381 & 1.0002 & 1.0000 & 1.0002 & 27.88 & 27.46 \\
039 & 708561 & 705615 & 703954 & 1.0042 & 1.0024 & 1.0065 & 37.02 & 39.21 \\
040 & 137746 & 137720 & 137487 & 1.0002 & 1.0017 & 1.0019 & 5.57 & 5.56 \\
041 & 2983392 & 2981318 & 2981318 & 1.0007 & 1.0000 & 1.0007 & 17.09 & 17.09 \\
042 & 844150 & 843446 & 843098 & 1.0008 & 1.0004 & 1.0012 & 37.09 & 37.08 \\
043 & 347613 & 347613 & 347613 & 1.0000 & 1.0000 & 1.0000 & 22.60 & 22.60 \\
044 & 68899 & 68899 & 67862 & 1.0000 & 1.0153 & 1.0153 & 0.06 & 0.06 \\
045 & 63059 & 63059 & 63052 & 1.0000 & 1.0001 & 1.0001 & 0.00 & 1.43 \\
046 & 166578 & 137552 & 125148 & 1.2110 & 1.0991 & 1.3310 & 57.36 & 67.58 \\
047 & 321355 & 321355 & 321355 & 1.0000 & 1.0000 & 1.0000 & 4.90 & 4.90 \\
048 & 219890 & 219890 & 219037 & 1.0000 & 1.0039 & 1.0039 & 4.18 & 4.31 \\
049 & 160206 & 160206 & 160206 & 1.0000 & 1.0000 & 1.0000 & 0.00 & 0.00 \\
050 & 108121 & 108121 & 107981 & 1.0000 & 1.0013 & 1.0013 & 1.07 & 3.52 \\
051 & 607628 & 607628 & 607628 & 1.0000 & 1.0000 & 1.0000 & 0.00 & 0.00 \\
052 & 90903 & 90903 & 90903 & 1.0000 & 1.0000 & 1.0000 & 11.27 & 11.27 \\
053 & 774929 & 768202 & 759346 & 1.0088 & 1.0117 & 1.0205 & 46.22 & 39.82 \\
054 & 517960 & 517960 & 513641 & 1.0000 & 1.0084 & 1.0084 & 0.00 & 0.00 \\
055 & 59839 & 59839 & 59839 & 1.0000 & 1.0000 & 1.0000 & 6.22 & 6.22 \\
056 & 265710 & 265710 & 265710 & 1.0000 & 1.0000 & 1.0000 & 0.00 & 0.00 \\
057 & 30183 & 30183 & 29658 & 1.0000 & 1.0177 & 1.0177 & 6.21 & 6.21 \\
058 & 2268376 & 2266610 & 2266610 & 1.0008 & 1.0000 & 1.0008 & 10.78 & 10.78 \\
059 & 369318 & 369318 & 369318 & 1.0000 & 1.0000 & 1.0000 & 0.00 & 0.00 \\
060 & 443955 & 443733 & 437319 & 1.0005 & 1.0147 & 1.0152 & 43.17 & 44.72 \\
061 & 39298 & 39272 & 38792 & 1.0007 & 1.0124 & 1.0130 & 5.97 & 7.26 \\
062 & 2521494 & 2521494 & 2387918 & 1.0000 & 1.0559 & 1.0559 & 0.00 & 0.00 \\
063 & 909394 & 903188 & 874911 & 1.0069 & 1.0323 & 1.0394 & 3.40 & 2.73 \\
064 & 19116 & 19116 & 19116 & 1.0000 & 1.0000 & 1.0000 & 0.00 & 0.00 \\
065 & 26610 & 26610 & 26551 & 1.0000 & 1.0022 & 1.0022 & 4.45 & 4.45 \\
066 & 248146 & 248146 & 247485 & 1.0000 & 1.0027 & 1.0027 & 0.00 & 0.00 \\
067 & 33768704 & 33768704 & 32164982 & 1.0000 & 1.0499 & 1.0499 & 0.00 & 7.19 \\
068 & 303697 & 303611 & 302692 & 1.0003 & 1.0030 & 1.0033 & 7.84 & 7.84 \\
069 & 25144 & 25144 & 24048 & 1.0000 & 1.0456 & 1.0456 & 0.00 & 0.00 \\
070 & 60121 & 60121 & 59140 & 1.0000 & 1.0166 & 1.0166 & 23.58 & 14.26 \\
071 & 17557 & 17557 & 17401 & 1.0000 & 1.0090 & 1.0090 & 0.00 & 0.00 \\
072 & 11833111 & 11819724 & 11819724 & 1.0011 & 1.0000 & 1.0011 & 8.67 & 8.67 \\
073 & 5980423 & 5980423 & 5176545 & 1.0000 & 1.1553 & 1.1553 & 0.00 & 8.85 \\
074 & 821757 & 820797 & 818681 & 1.0012 & 1.0026 & 1.0038 & 23.89 & 24.00 \\
075 & 1397410 & 1397410 & 1389088 & 1.0000 & 1.0060 & 1.0060 & 5.63 & 5.63 \\
076 & 3065214 & 3063843 & 3063504 & 1.0004 & 1.0001 & 1.0006 & 27.88 & 27.47 \\
077 & 248734 & 248734 & 248734 & 1.0000 & 1.0000 & 1.0000 & 2.61 & 2.61 \\
078 & 1260358 & 1260358 & 1149984 & 1.0000 & 1.0960 & 1.0960 & 0.00 & 23.16 \\
079 & 15446486 & 15446486 & 15446486 & 1.0000 & 1.0000 & 1.0000 & 0.00 & 0.00 \\
080 & 145926 & 145674 & 145674 & 1.0017 & 1.0000 & 1.0017 & 43.60 & 43.60 \\
081 & 76860 & 76860 & 76860 & 1.0000 & 1.0000 & 1.0000 & 0.00 & 0.00 \\
082 & 708711 & 708711 & 708711 & 1.0000 & 1.0000 & 1.0000 & 7.65 & 7.65 \\
083 & 644672 & 535329 & 465958 & 1.2043 & 1.1489 & 1.3835 & 70.76 & 68.77 \\
084 & 1255052 & 1255052 & 1255052 & 1.0000 & 1.0000 & 1.0000 & 0.00 & 0.00 \\
085 & 3168404 & 3168404 & 3157230 & 1.0000 & 1.0035 & 1.0035 & 8.22 & 8.22 \\
086 & 118403 & 118403 & 115496 & 1.0000 & 1.0252 & 1.0252 & 0.00 & 0.00 \\
087 & 1984436 & 1982828 & 1982828 & 1.0008 & 1.0000 & 1.0008 & 26.83 & 26.83 \\
088 & 203492 & 203492 & 203492 & 1.0000 & 1.0000 & 1.0000 & 16.61 & 16.61 \\
089 & 346832 & 346832 & 346763 & 1.0000 & 1.0002 & 1.0002 & 0.00 & 35.08 \\
090 & 677936 & 599940 & 549431 & 1.1300 & 1.0919 & 1.2339 & 41.08 & 65.00 \\
091 & 4571759 & 4569395 & 4569395 & 1.0005 & 1.0000 & 1.0005 & 17.99 & 17.99 \\
092 & 2828252 & 2358868 & 2358868 & 1.1990 & 1.0000 & 1.1990 & 75.92 & 75.92 \\
093 & 37827 & 37827 & 37827 & 1.0000 & 1.0000 & 1.0000 & 0.00 & 0.00 \\
094 & 48073 & 48073 & 46927 & 1.0000 & 1.0244 & 1.0244 & 6.43 & 18.40 \\
095 & 1042890 & 1042890 & 1042890 & 1.0000 & 1.0000 & 1.0000 & 0.00 & 0.00 \\
096 & 44010 & 44010 & 41802 & 1.0000 & 1.0528 & 1.0528 & 24.27 & 24.27 \\
097 & 5617602 & 5617602 & 5617602 & 1.0000 & 1.0000 & 1.0000 & 0.00 & 0.00 \\
098 & 89756 & 89756 & 89756 & 1.0000 & 1.0000 & 1.0000 & 0.00 & 0.00 \\
099 & 808916 & 806053 & 806053 & 1.0036 & 1.0000 & 1.0036 & 43.77 & 43.77 \\
100 & 81080 & 81080 & 77421 & 1.0000 & 1.0473 & 1.0473 & 3.73 & 3.73 \\
\end{longtable}

\begin{longtable}{crrrrrrrr}
\caption{3核Cache配对（100图）}\label{tab:cache-detail-3}\\
\toprule
图 & $T_B$ & $T_C(X_B)$ & $T_C(X_C)$ & $S_{hw}$ & $S_{ad}$ & $S_{opt}$ & $C(X_B)$命中/\% & $C(X_C)$命中/\% \\
\midrule
\endfirsthead
\multicolumn{9}{c}{续表}\\
\toprule
图 & $T_B$ & $T_C(X_B)$ & $T_C(X_C)$ & $S_{hw}$ & $S_{ad}$ & $S_{opt}$ & $C(X_B)$命中/\% & $C(X_C)$命中/\% \\
\midrule
\endhead
\midrule\multicolumn{9}{r}{续下页}\\
\endfoot
\bottomrule
\endlastfoot
001 & 78545 & 78545 & 78545 & 1.0000 & 1.0000 & 1.0000 & 0.00 & 0.00 \\
002 & 131349 & 131327 & 127539 & 1.0002 & 1.0297 & 1.0299 & 15.92 & 15.66 \\
003 & 524210 & 524167 & 522594 & 1.0001 & 1.0030 & 1.0031 & 11.83 & 10.75 \\
004 & 134043 & 134043 & 134043 & 1.0000 & 1.0000 & 1.0000 & 0.00 & 0.00 \\
005 & 88696 & 88666 & 86911 & 1.0003 & 1.0202 & 1.0205 & 11.84 & 11.84 \\
006 & 31152 & 31152 & 31152 & 1.0000 & 1.0000 & 1.0000 & 5.25 & 5.25 \\
007 & 338094 & 338094 & 338094 & 1.0000 & 1.0000 & 1.0000 & 0.00 & 0.00 \\
008 & 163359 & 163359 & 163359 & 1.0000 & 1.0000 & 1.0000 & 0.00 & 0.00 \\
009 & 91278 & 91278 & 91278 & 1.0000 & 1.0000 & 1.0000 & 0.01 & 0.01 \\
010 & 32786 & 32786 & 30412 & 1.0000 & 1.0781 & 1.0781 & 4.64 & 9.94 \\
011 & 69279 & 69144 & 68637 & 1.0020 & 1.0074 & 1.0094 & 60.38 & 60.71 \\
012 & 20661 & 20661 & 20561 & 1.0000 & 1.0049 & 1.0049 & 7.49 & 7.47 \\
013 & 88351 & 88351 & 88351 & 1.0000 & 1.0000 & 1.0000 & 0.00 & 0.00 \\
014 & 5980663 & 5942569 & 5942569 & 1.0064 & 1.0000 & 1.0064 & 28.89 & 28.89 \\
015 & 60578 & 60578 & 60578 & 1.0000 & 1.0000 & 1.0000 & 24.69 & 24.69 \\
016 & 7712915 & 7712915 & 7712915 & 1.0000 & 1.0000 & 1.0000 & 0.00 & 0.00 \\
017 & 55389 & 55389 & 55154 & 1.0000 & 1.0043 & 1.0043 & 12.98 & 13.73 \\
018 & 341047 & 338488 & 338488 & 1.0076 & 1.0000 & 1.0076 & 62.56 & 62.56 \\
019 & 28657 & 28657 & 26572 & 1.0000 & 1.0785 & 1.0785 & 23.26 & 23.26 \\
020 & 177950 & 177950 & 177950 & 1.0000 & 1.0000 & 1.0000 & 0.00 & 0.00 \\
021 & 2144024 & 2142851 & 2142612 & 1.0005 & 1.0001 & 1.0007 & 54.59 & 56.35 \\
022 & 19026 & 19026 & 18965 & 1.0000 & 1.0032 & 1.0032 & 6.08 & 6.08 \\
023 & 40897 & 40897 & 40897 & 1.0000 & 1.0000 & 1.0000 & 32.35 & 32.35 \\
024 & 2552735 & 2552735 & 2552735 & 1.0000 & 1.0000 & 1.0000 & 0.00 & 0.00 \\
025 & 1622164 & 1622164 & 1622138 & 1.0000 & 1.0000 & 1.0000 & 0.00 & 0.70 \\
026 & 41665 & 41665 & 41665 & 1.0000 & 1.0000 & 1.0000 & 16.37 & 16.37 \\
027 & 831811 & 831458 & 831365 & 1.0004 & 1.0001 & 1.0005 & 50.64 & 51.08 \\
028 & 16770711 & 16770711 & 16770711 & 1.0000 & 1.0000 & 1.0000 & 0.00 & 0.00 \\
029 & 12743 & 12743 & 12743 & 1.0000 & 1.0000 & 1.0000 & 0.00 & 0.00 \\
030 & 933857 & 933236 & 932731 & 1.0007 & 1.0005 & 1.0012 & 54.65 & 54.63 \\
031 & 473822 & 473000 & 473000 & 1.0017 & 1.0000 & 1.0017 & 49.73 & 49.73 \\
032 & 14175 & 14175 & 14175 & 1.0000 & 1.0000 & 1.0000 & 4.26 & 4.26 \\
033 & 592065 & 590152 & 590152 & 1.0032 & 1.0000 & 1.0032 & 38.26 & 38.26 \\
034 & 276994 & 276560 & 276560 & 1.0016 & 1.0000 & 1.0016 & 54.75 & 54.75 \\
035 & 95406 & 95406 & 95406 & 1.0000 & 1.0000 & 1.0000 & 0.00 & 0.00 \\
036 & 311345 & 311345 & 311345 & 1.0000 & 1.0000 & 1.0000 & 0.00 & 0.00 \\
037 & 74038 & 74038 & 74038 & 1.0000 & 1.0000 & 1.0000 & 58.97 & 58.97 \\
038 & 452248 & 452022 & 452022 & 1.0005 & 1.0000 & 1.0005 & 35.73 & 35.73 \\
039 & 485565 & 483226 & 482066 & 1.0048 & 1.0024 & 1.0073 & 38.60 & 39.33 \\
040 & 140837 & 140811 & 140543 & 1.0002 & 1.0019 & 1.0021 & 1.33 & 1.33 \\
041 & 1991911 & 1989953 & 1989953 & 1.0010 & 1.0000 & 1.0010 & 26.05 & 26.05 \\
042 & 562868 & 561628 & 561628 & 1.0022 & 1.0000 & 1.0022 & 43.33 & 43.33 \\
043 & 268777 & 268777 & 265478 & 1.0000 & 1.0124 & 1.0124 & 8.99 & 8.99 \\
044 & 68899 & 68899 & 67862 & 1.0000 & 1.0153 & 1.0153 & 0.06 & 0.06 \\
045 & 42601 & 42601 & 42601 & 1.0000 & 1.0000 & 1.0000 & 0.00 & 0.00 \\
046 & 142176 & 121494 & 106079 & 1.1702 & 1.1453 & 1.3403 & 43.06 & 51.09 \\
047 & 326508 & 326508 & 324264 & 1.0000 & 1.0069 & 1.0069 & 0.00 & 0.00 \\
048 & 218425 & 218419 & 218419 & 1.0000 & 1.0000 & 1.0000 & 8.90 & 8.90 \\
049 & 160206 & 160206 & 160206 & 1.0000 & 1.0000 & 1.0000 & 0.00 & 0.00 \\
050 & 108121 & 108121 & 107981 & 1.0000 & 1.0013 & 1.0013 & 1.07 & 3.52 \\
051 & 607628 & 607628 & 607628 & 1.0000 & 1.0000 & 1.0000 & 0.00 & 0.00 \\
052 & 61071 & 61071 & 61071 & 1.0000 & 1.0000 & 1.0000 & 18.28 & 18.28 \\
053 & 542159 & 506885 & 491765 & 1.0696 & 1.0307 & 1.1025 & 37.07 & 43.81 \\
054 & 516442 & 516442 & 513130 & 1.0000 & 1.0065 & 1.0065 & 0.00 & 0.00 \\
055 & 42047 & 42047 & 39987 & 1.0000 & 1.0515 & 1.0515 & 15.61 & 14.50 \\
056 & 265710 & 265710 & 265710 & 1.0000 & 1.0000 & 1.0000 & 0.00 & 0.00 \\
057 & 20705 & 20705 & 20705 & 1.0000 & 1.0000 & 1.0000 & 0.08 & 0.08 \\
058 & 1554831 & 1554831 & 1533246 & 1.0000 & 1.0141 & 1.0141 & 0.00 & 17.50 \\
059 & 266217 & 262333 & 257883 & 1.0148 & 1.0173 & 1.0323 & 68.32 & 71.87 \\
060 & 298210 & 297398 & 297002 & 1.0027 & 1.0013 & 1.0041 & 48.04 & 47.86 \\
061 & 30303 & 30303 & 30303 & 1.0000 & 1.0000 & 1.0000 & 6.22 & 6.22 \\
062 & 1939281 & 1939281 & 1818680 & 1.0000 & 1.0663 & 1.0663 & 0.00 & 0.00 \\
063 & 617304 & 616709 & 593484 & 1.0010 & 1.0391 & 1.0401 & 0.27 & 0.27 \\
064 & 19116 & 19116 & 19116 & 1.0000 & 1.0000 & 1.0000 & 0.00 & 0.00 \\
065 & 17972 & 17972 & 17972 & 1.0000 & 1.0000 & 1.0000 & 8.53 & 8.53 \\
066 & 248146 & 248146 & 247485 & 1.0000 & 1.0027 & 1.0027 & 0.00 & 0.00 \\
067 & 23568859 & 23568859 & 22573181 & 1.0000 & 1.0441 & 1.0441 & 0.00 & 5.99 \\
068 & 281918 & 281418 & 281418 & 1.0018 & 1.0000 & 1.0018 & 17.23 & 17.23 \\
069 & 24263 & 24246 & 22922 & 1.0007 & 1.0578 & 1.0585 & 8.01 & 8.26 \\
070 & 39333 & 39333 & 39333 & 1.0000 & 1.0000 & 1.0000 & 27.30 & 27.30 \\
071 & 17557 & 17557 & 17401 & 1.0000 & 1.0090 & 1.0090 & 0.00 & 0.00 \\
072 & 7935290 & 7922729 & 7921555 & 1.0016 & 1.0001 & 1.0017 & 15.41 & 15.44 \\
073 & 4515227 & 4515227 & 3969025 & 1.0000 & 1.1376 & 1.1376 & 0.00 & 10.66 \\
074 & 559584 & 555925 & 554482 & 1.0066 & 1.0026 & 1.0092 & 32.40 & 32.22 \\
075 & 1363452 & 1363444 & 1363444 & 1.0000 & 1.0000 & 1.0000 & 10.09 & 10.09 \\
076 & 2050028 & 2047868 & 2047868 & 1.0011 & 1.0000 & 1.0011 & 33.69 & 33.69 \\
077 & 168976 & 168976 & 168126 & 1.0000 & 1.0051 & 1.0051 & 2.54 & 2.54 \\
078 & 993959 & 938948 & 882747 & 1.0586 & 1.0637 & 1.1260 & 22.46 & 33.34 \\
079 & 10181692 & 10175315 & 10175315 & 1.0006 & 1.0000 & 1.0006 & 13.73 & 13.73 \\
080 & 100325 & 100325 & 100325 & 1.0000 & 1.0000 & 1.0000 & 0.00 & 0.00 \\
081 & 51943 & 51943 & 51943 & 1.0000 & 1.0000 & 1.0000 & 0.00 & 0.00 \\
082 & 619074 & 618927 & 606773 & 1.0002 & 1.0200 & 1.0203 & 7.44 & 7.60 \\
083 & 512623 & 402691 & 402691 & 1.2730 & 1.0000 & 1.2730 & 65.90 & 65.90 \\
084 & 836220 & 836220 & 836220 & 1.0000 & 1.0000 & 1.0000 & 0.00 & 0.00 \\
085 & 3130505 & 3130505 & 3112095 & 1.0000 & 1.0059 & 1.0059 & 0.00 & 0.00 \\
086 & 118403 & 118403 & 115496 & 1.0000 & 1.0252 & 1.0252 & 0.00 & 0.00 \\
087 & 1374901 & 1372451 & 1372451 & 1.0018 & 1.0000 & 1.0018 & 26.15 & 26.15 \\
088 & 196677 & 195998 & 190688 & 1.0035 & 1.0278 & 1.0314 & 20.78 & 20.78 \\
089 & 243614 & 243614 & 242586 & 1.0000 & 1.0042 & 1.0042 & 0.00 & 25.56 \\
090 & 657430 & 653823 & 549742 & 1.0055 & 1.1893 & 1.1959 & 0.39 & 19.33 \\
091 & 3078087 & 3057457 & 3057457 & 1.0067 & 1.0000 & 1.0067 & 23.43 & 23.43 \\
092 & 2231804 & 1792534 & 1792534 & 1.2451 & 1.0000 & 1.2451 & 65.42 & 65.42 \\
093 & 25103 & 25103 & 25103 & 1.0000 & 1.0000 & 1.0000 & 0.00 & 0.00 \\
094 & 32222 & 32222 & 31679 & 1.0000 & 1.0171 & 1.0171 & 31.08 & 28.38 \\
095 & 695775 & 695775 & 695775 & 1.0000 & 1.0000 & 1.0000 & 0.00 & 0.00 \\
096 & 30318 & 30318 & 30318 & 1.0000 & 1.0000 & 1.0000 & 27.43 & 27.43 \\
097 & 3660459 & 3660338 & 3622881 & 1.0000 & 1.0103 & 1.0104 & 21.97 & 25.99 \\
098 & 89756 & 89756 & 89756 & 1.0000 & 1.0000 & 1.0000 & 0.00 & 0.00 \\
099 & 539533 & 538112 & 537694 & 1.0026 & 1.0008 & 1.0034 & 43.31 & 43.39 \\
100 & 57811 & 57811 & 57811 & 1.0000 & 1.0000 & 1.0000 & 10.58 & 10.58 \\
\end{longtable}

\begin{longtable}{crrrrrrrr}
\caption{4核Cache配对（100图）}\label{tab:cache-detail-4}\\
\toprule
图 & $T_B$ & $T_C(X_B)$ & $T_C(X_C)$ & $S_{hw}$ & $S_{ad}$ & $S_{opt}$ & $C(X_B)$命中/\% & $C(X_C)$命中/\% \\
\midrule
\endfirsthead
\multicolumn{9}{c}{续表}\\
\toprule
图 & $T_B$ & $T_C(X_B)$ & $T_C(X_C)$ & $S_{hw}$ & $S_{ad}$ & $S_{opt}$ & $C(X_B)$命中/\% & $C(X_C)$命中/\% \\
\midrule
\endhead
\midrule\multicolumn{9}{r}{续下页}\\
\endfoot
\bottomrule
\endlastfoot
001 & 58984 & 58984 & 58984 & 1.0000 & 1.0000 & 1.0000 & 0.00 & 0.00 \\
002 & 122326 & 120380 & 119922 & 1.0162 & 1.0038 & 1.0200 & 10.96 & 10.71 \\
003 & 524210 & 524167 & 522594 & 1.0001 & 1.0030 & 1.0031 & 11.83 & 10.75 \\
004 & 104188 & 104188 & 104188 & 1.0000 & 1.0000 & 1.0000 & 0.00 & 0.00 \\
005 & 88696 & 88666 & 86911 & 1.0003 & 1.0202 & 1.0205 & 11.84 & 11.84 \\
006 & 31152 & 31152 & 31152 & 1.0000 & 1.0000 & 1.0000 & 5.25 & 5.25 \\
007 & 254398 & 254398 & 254398 & 1.0000 & 1.0000 & 1.0000 & 0.00 & 0.00 \\
008 & 123060 & 123060 & 123060 & 1.0000 & 1.0000 & 1.0000 & 0.00 & 0.00 \\
009 & 91278 & 91278 & 90556 & 1.0000 & 1.0080 & 1.0080 & 0.01 & 5.55 \\
010 & 29918 & 29858 & 29845 & 1.0020 & 1.0004 & 1.0024 & 10.30 & 13.66 \\
011 & 46880 & 46880 & 46880 & 1.0000 & 1.0000 & 1.0000 & 0.00 & 0.00 \\
012 & 15863 & 15863 & 15705 & 1.0000 & 1.0101 & 1.0101 & 9.28 & 9.28 \\
013 & 66982 & 66982 & 66982 & 1.0000 & 1.0000 & 1.0000 & 0.00 & 0.00 \\
014 & 4521815 & 4500714 & 4500714 & 1.0047 & 1.0000 & 1.0047 & 27.41 & 27.41 \\
015 & 48727 & 48669 & 48669 & 1.0012 & 1.0000 & 1.0012 & 20.89 & 20.89 \\
016 & 7706495 & 7706495 & 7706472 & 1.0000 & 1.0000 & 1.0000 & 12.62 & 12.74 \\
017 & 41209 & 41209 & 41209 & 1.0000 & 1.0000 & 1.0000 & 31.85 & 31.85 \\
018 & 253676 & 253384 & 253384 & 1.0012 & 1.0000 & 1.0012 & 56.24 & 56.24 \\
019 & 24099 & 24099 & 23835 & 1.0000 & 1.0111 & 1.0111 & 19.75 & 20.09 \\
020 & 133704 & 133704 & 133704 & 1.0000 & 1.0000 & 1.0000 & 0.00 & 0.00 \\
021 & 1550324 & 1550324 & 1550324 & 1.0000 & 1.0000 & 1.0000 & 0.00 & 0.00 \\
022 & 14526 & 14526 & 14275 & 1.0000 & 1.0176 & 1.0176 & 6.08 & 6.08 \\
023 & 30142 & 30142 & 30142 & 1.0000 & 1.0000 & 1.0000 & 39.32 & 39.32 \\
024 & 2552735 & 2552735 & 2552735 & 1.0000 & 1.0000 & 1.0000 & 0.00 & 0.00 \\
025 & 1217444 & 1217444 & 1217444 & 1.0000 & 1.0000 & 1.0000 & 0.00 & 0.00 \\
026 & 41665 & 41665 & 41665 & 1.0000 & 1.0000 & 1.0000 & 16.37 & 16.37 \\
027 & 559196 & 559196 & 559196 & 1.0000 & 1.0000 & 1.0000 & 0.00 & 0.00 \\
028 & 13608312 & 13585418 & 13585418 & 1.0017 & 1.0000 & 1.0017 & 0.27 & 0.27 \\
029 & 11028 & 11028 & 11028 & 1.0000 & 1.0000 & 1.0000 & 0.00 & 0.00 \\
030 & 700242 & 699694 & 699694 & 1.0008 & 1.0000 & 1.0008 & 52.45 & 52.45 \\
031 & 357987 & 355839 & 355839 & 1.0060 & 1.0000 & 1.0060 & 48.80 & 48.80 \\
032 & 11004 & 11004 & 11004 & 1.0000 & 1.0000 & 1.0000 & 6.24 & 6.24 \\
033 & 443447 & 442667 & 441850 & 1.0018 & 1.0018 & 1.0036 & 42.67 & 41.73 \\
034 & 212170 & 212170 & 212170 & 1.0000 & 1.0000 & 1.0000 & 52.33 & 52.33 \\
035 & 95406 & 95406 & 95406 & 1.0000 & 1.0000 & 1.0000 & 0.00 & 0.00 \\
036 & 233584 & 233584 & 233584 & 1.0000 & 1.0000 & 1.0000 & 0.00 & 0.00 \\
037 & 53320 & 53320 & 53320 & 1.0000 & 1.0000 & 1.0000 & 0.00 & 0.00 \\
038 & 339624 & 337813 & 337813 & 1.0054 & 1.0000 & 1.0054 & 35.21 & 35.21 \\
039 & 395323 & 368055 & 368055 & 1.0741 & 1.0000 & 1.0741 & 47.82 & 47.82 \\
040 & 140837 & 140811 & 136073 & 1.0002 & 1.0348 & 1.0350 & 1.33 & 2.28 \\
041 & 1497746 & 1494017 & 1494017 & 1.0025 & 1.0000 & 1.0025 & 24.13 & 24.13 \\
042 & 426511 & 423652 & 423652 & 1.0067 & 1.0000 & 1.0067 & 43.14 & 43.14 \\
043 & 295497 & 295497 & 295497 & 1.0000 & 1.0000 & 1.0000 & 18.70 & 18.70 \\
044 & 68899 & 68899 & 67862 & 1.0000 & 1.0153 & 1.0153 & 0.06 & 0.06 \\
045 & 31878 & 31878 & 31878 & 1.0000 & 1.0000 & 1.0000 & 0.00 & 0.00 \\
046 & 124612 & 105932 & 98844 & 1.1763 & 1.0717 & 1.2607 & 31.27 & 32.38 \\
047 & 314757 & 314754 & 314754 & 1.0000 & 1.0000 & 1.0000 & 10.25 & 10.25 \\
048 & 218425 & 218419 & 218419 & 1.0000 & 1.0000 & 1.0000 & 8.90 & 8.90 \\
049 & 160206 & 160206 & 160206 & 1.0000 & 1.0000 & 1.0000 & 0.00 & 0.00 \\
050 & 107794 & 107794 & 107658 & 1.0000 & 1.0013 & 1.0013 & 1.04 & 1.04 \\
051 & 541017 & 541017 & 541017 & 1.0000 & 1.0000 & 1.0000 & 8.33 & 8.33 \\
052 & 45610 & 45610 & 45610 & 1.0000 & 1.0000 & 1.0000 & 16.59 & 16.59 \\
053 & 492232 & 492232 & 491765 & 1.0000 & 1.0009 & 1.0009 & 35.50 & 43.81 \\
054 & 516442 & 516442 & 513130 & 1.0000 & 1.0065 & 1.0065 & 0.00 & 0.00 \\
055 & 32099 & 32099 & 32099 & 1.0000 & 1.0000 & 1.0000 & 21.72 & 21.72 \\
056 & 265710 & 265710 & 265710 & 1.0000 & 1.0000 & 1.0000 & 0.00 & 0.00 \\
057 & 16485 & 16485 & 16485 & 1.0000 & 1.0000 & 1.0000 & 7.74 & 7.74 \\
058 & 1154519 & 1142629 & 1142629 & 1.0104 & 1.0000 & 1.0104 & 26.90 & 26.90 \\
059 & 185216 & 185216 & 185216 & 1.0000 & 1.0000 & 1.0000 & 0.00 & 0.00 \\
060 & 221720 & 221694 & 221694 & 1.0001 & 1.0000 & 1.0001 & 48.17 & 47.84 \\
061 & 24024 & 24024 & 24024 & 1.0000 & 1.0000 & 1.0000 & 11.49 & 11.49 \\
062 & 1767455 & 1270730 & 1270730 & 1.3909 & 1.0000 & 1.3909 & 0.03 & 0.03 \\
063 & 324251 & 324251 & 324251 & 1.0000 & 1.0000 & 1.0000 & 0.00 & 0.00 \\
064 & 19116 & 19116 & 19116 & 1.0000 & 1.0000 & 1.0000 & 0.00 & 0.00 \\
065 & 13244 & 13244 & 13062 & 1.0000 & 1.0139 & 1.0139 & 8.18 & 9.53 \\
066 & 248146 & 248146 & 247485 & 1.0000 & 1.0027 & 1.0027 & 0.00 & 0.00 \\
067 & 18473155 & 18473155 & 17693167 & 1.0000 & 1.0441 & 1.0441 & 0.00 & 7.27 \\
068 & 267505 & 266974 & 266974 & 1.0020 & 1.0000 & 1.0020 & 20.71 & 20.71 \\
069 & 24304 & 24134 & 24134 & 1.0070 & 1.0000 & 1.0070 & 10.84 & 10.84 \\
070 & 31225 & 31192 & 31192 & 1.0011 & 1.0000 & 1.0011 & 39.93 & 39.93 \\
071 & 17557 & 17557 & 17401 & 1.0000 & 1.0090 & 1.0090 & 0.00 & 0.00 \\
072 & 5993260 & 5988828 & 5988828 & 1.0007 & 1.0000 & 1.0007 & 10.69 & 10.69 \\
073 & 3903787 & 3903787 & 3434743 & 1.0000 & 1.1366 & 1.1366 & 0.00 & 11.05 \\
074 & 418987 & 415633 & 415633 & 1.0081 & 1.0000 & 1.0081 & 32.23 & 32.23 \\
075 & 1320997 & 1320971 & 1320971 & 1.0000 & 1.0000 & 1.0000 & 15.06 & 15.06 \\
076 & 1540201 & 1536306 & 1536296 & 1.0025 & 1.0000 & 1.0025 & 41.71 & 41.67 \\
077 & 127832 & 127591 & 126652 & 1.0019 & 1.0074 & 1.0093 & 5.08 & 5.07 \\
078 & 690924 & 690924 & 642020 & 1.0000 & 1.0762 & 1.0762 & 0.00 & 10.99 \\
079 & 7635624 & 7635624 & 7635624 & 1.0000 & 1.0000 & 1.0000 & 0.00 & 0.00 \\
080 & 111314 & 90584 & 90559 & 1.2288 & 1.0003 & 1.2292 & 47.85 & 49.02 \\
081 & 39588 & 39588 & 39588 & 1.0000 & 1.0000 & 1.0000 & 0.00 & 0.00 \\
082 & 581815 & 581598 & 581598 & 1.0004 & 1.0000 & 1.0004 & 25.22 & 25.22 \\
083 & 450029 & 333151 & 333151 & 1.3508 & 1.0000 & 1.3508 & 64.26 & 64.26 \\
084 & 629028 & 629028 & 629028 & 1.0000 & 1.0000 & 1.0000 & 0.00 & 0.00 \\
085 & 2987335 & 2986917 & 2986917 & 1.0001 & 1.0000 & 1.0001 & 16.23 & 16.23 \\
086 & 94047 & 94047 & 88995 & 1.0000 & 1.0568 & 1.0568 & 4.51 & 4.51 \\
087 & 1046846 & 1043492 & 1042446 & 1.0032 & 1.0010 & 1.0042 & 29.95 & 29.97 \\
088 & 185641 & 184665 & 184608 & 1.0053 & 1.0003 & 1.0056 & 22.34 & 23.19 \\
089 & 178160 & 178160 & 175147 & 1.0000 & 1.0172 & 1.0172 & 0.00 & 20.73 \\
090 & 608275 & 604430 & 534087 & 1.0064 & 1.1317 & 1.1389 & 0.45 & 6.81 \\
091 & 2314405 & 2306119 & 2306119 & 1.0036 & 1.0000 & 1.0036 & 19.88 & 19.88 \\
092 & 1927728 & 1470508 & 1470508 & 1.3109 & 1.0000 & 1.3109 & 64.79 & 64.79 \\
093 & 19779 & 19779 & 19779 & 1.0000 & 1.0000 & 1.0000 & 0.00 & 0.00 \\
094 & 24640 & 24640 & 24640 & 1.0000 & 1.0000 & 1.0000 & 38.58 & 38.58 \\
095 & 524422 & 524422 & 524422 & 1.0000 & 1.0000 & 1.0000 & 0.00 & 0.00 \\
096 & 26984 & 26984 & 26984 & 1.0000 & 1.0000 & 1.0000 & 8.35 & 8.35 \\
097 & 2626908 & 2626908 & 2626908 & 1.0000 & 1.0000 & 1.0000 & 4.17 & 4.17 \\
098 & 89756 & 89756 & 89756 & 1.0000 & 1.0000 & 1.0000 & 0.00 & 0.00 \\
099 & 405336 & 403555 & 403147 & 1.0044 & 1.0010 & 1.0054 & 55.68 & 53.60 \\
100 & 87316 & 87316 & 72103 & 1.0000 & 1.2110 & 1.2110 & 2.74 & 25.82 \\
\end{longtable}

\begin{longtable}{crrrrrrrr}
\caption{5核Cache配对（100图）}\label{tab:cache-detail-5}\\
\toprule
图 & $T_B$ & $T_C(X_B)$ & $T_C(X_C)$ & $S_{hw}$ & $S_{ad}$ & $S_{opt}$ & $C(X_B)$命中/\% & $C(X_C)$命中/\% \\
\midrule
\endfirsthead
\multicolumn{9}{c}{续表}\\
\toprule
图 & $T_B$ & $T_C(X_B)$ & $T_C(X_C)$ & $S_{hw}$ & $S_{ad}$ & $S_{opt}$ & $C(X_B)$命中/\% & $C(X_C)$命中/\% \\
\midrule
\endhead
\midrule\multicolumn{9}{r}{续下页}\\
\endfoot
\bottomrule
\endlastfoot
001 & 47502 & 47502 & 47502 & 1.0000 & 1.0000 & 1.0000 & 0.00 & 0.00 \\
002 & 111678 & 108499 & 105055 & 1.0293 & 1.0328 & 1.0630 & 10.18 & 10.84 \\
003 & 469702 & 469516 & 469516 & 1.0004 & 1.0000 & 1.0004 & 18.29 & 18.29 \\
004 & 88341 & 88341 & 86005 & 1.0000 & 1.0272 & 1.0272 & 0.00 & 20.88 \\
005 & 88696 & 88666 & 86911 & 1.0003 & 1.0202 & 1.0205 & 11.84 & 11.84 \\
006 & 31152 & 31152 & 31152 & 1.0000 & 1.0000 & 1.0000 & 5.25 & 5.25 \\
007 & 203834 & 203834 & 203834 & 1.0000 & 1.0000 & 1.0000 & 0.00 & 0.00 \\
008 & 100603 & 100603 & 100603 & 1.0000 & 1.0000 & 1.0000 & 0.00 & 0.00 \\
009 & 91278 & 91278 & 90556 & 1.0000 & 1.0080 & 1.0080 & 0.01 & 5.55 \\
010 & 29918 & 29858 & 29845 & 1.0020 & 1.0004 & 1.0024 & 10.30 & 13.66 \\
011 & 46880 & 46880 & 46777 & 1.0000 & 1.0022 & 1.0022 & 0.00 & 62.31 \\
012 & 13406 & 13406 & 13406 & 1.0000 & 1.0000 & 1.0000 & 7.68 & 7.68 \\
013 & 54349 & 54125 & 53418 & 1.0041 & 1.0132 & 1.0174 & 1.19 & 1.49 \\
014 & 3681597 & 3672221 & 3672221 & 1.0026 & 1.0000 & 1.0026 & 27.00 & 27.00 \\
015 & 43665 & 42872 & 42872 & 1.0185 & 1.0000 & 1.0185 & 24.87 & 24.87 \\
016 & 7706495 & 7706495 & 7706472 & 1.0000 & 1.0000 & 1.0000 & 12.62 & 12.74 \\
017 & 34231 & 34231 & 34231 & 1.0000 & 1.0000 & 1.0000 & 21.62 & 21.62 \\
018 & 205219 & 201659 & 201659 & 1.0177 & 1.0000 & 1.0177 & 57.22 & 57.22 \\
019 & 24099 & 24099 & 23835 & 1.0000 & 1.0111 & 1.0111 & 19.75 & 20.09 \\
020 & 107650 & 107650 & 107650 & 1.0000 & 1.0000 & 1.0000 & 0.00 & 0.00 \\
021 & 1250168 & 1249380 & 1249380 & 1.0006 & 1.0000 & 1.0006 & 57.00 & 57.00 \\
022 & 11824 & 11824 & 11824 & 1.0000 & 1.0000 & 1.0000 & 8.85 & 8.85 \\
023 & 24517 & 24517 & 24517 & 1.0000 & 1.0000 & 1.0000 & 21.00 & 21.00 \\
024 & 2538844 & 2538844 & 2538844 & 1.0000 & 1.0000 & 1.0000 & 18.94 & 18.94 \\
025 & 974747 & 974747 & 974747 & 1.0000 & 1.0000 & 1.0000 & 0.00 & 0.00 \\
026 & 41665 & 41665 & 41665 & 1.0000 & 1.0000 & 1.0000 & 16.37 & 16.37 \\
027 & 557147 & 556819 & 552561 & 1.0006 & 1.0077 & 1.0083 & 11.79 & 12.21 \\
028 & 11491035 & 11450794 & 10636977 & 1.0035 & 1.0765 & 1.0803 & 1.06 & 14.99 \\
029 & 9383 & 9383 & 9383 & 1.0000 & 1.0000 & 1.0000 & 0.00 & 0.00 \\
030 & 559889 & 559872 & 559025 & 1.0000 & 1.0015 & 1.0015 & 54.77 & 54.77 \\
031 & 285693 & 284489 & 284489 & 1.0042 & 1.0000 & 1.0042 & 53.15 & 53.15 \\
032 & 10496 & 10479 & 10479 & 1.0016 & 1.0000 & 1.0016 & 4.26 & 4.26 \\
033 & 355102 & 354514 & 354514 & 1.0017 & 1.0000 & 1.0017 & 31.27 & 31.27 \\
034 & 170371 & 168925 & 168887 & 1.0086 & 1.0002 & 1.0088 & 58.25 & 58.62 \\
035 & 95406 & 95406 & 95406 & 1.0000 & 1.0000 & 1.0000 & 0.00 & 0.00 \\
036 & 187182 & 187182 & 187182 & 1.0000 & 1.0000 & 1.0000 & 0.00 & 0.00 \\
037 & 43800 & 43013 & 43013 & 1.0183 & 1.0000 & 1.0183 & 64.25 & 64.25 \\
038 & 275133 & 272533 & 272533 & 1.0095 & 1.0000 & 1.0095 & 42.22 & 42.22 \\
039 & 395323 & 368055 & 340476 & 1.0741 & 1.0810 & 1.1611 & 47.82 & 37.09 \\
040 & 140837 & 140811 & 140543 & 1.0002 & 1.0019 & 1.0021 & 1.33 & 1.33 \\
041 & 1201884 & 1194643 & 1194643 & 1.0061 & 1.0000 & 1.0061 & 32.66 & 32.66 \\
042 & 342520 & 339457 & 339457 & 1.0090 & 1.0000 & 1.0090 & 49.37 & 49.37 \\
043 & 295497 & 295497 & 295497 & 1.0000 & 1.0000 & 1.0000 & 18.70 & 18.70 \\
044 & 68899 & 68899 & 63439 & 1.0000 & 1.0861 & 1.0861 & 0.06 & 0.00 \\
045 & 26164 & 26164 & 26164 & 1.0000 & 1.0000 & 1.0000 & 0.00 & 0.00 \\
046 & 117204 & 103554 & 98844 & 1.1318 & 1.0477 & 1.1857 & 23.05 & 32.38 \\
047 & 316261 & 316261 & 315056 & 1.0000 & 1.0038 & 1.0038 & 0.05 & 0.05 \\
048 & 217218 & 217140 & 215393 & 1.0004 & 1.0081 & 1.0085 & 11.32 & 11.32 \\
049 & 160206 & 160206 & 160206 & 1.0000 & 1.0000 & 1.0000 & 0.00 & 0.00 \\
050 & 107794 & 107794 & 107658 & 1.0000 & 1.0013 & 1.0013 & 1.04 & 1.04 \\
051 & 541017 & 541017 & 541017 & 1.0000 & 1.0000 & 1.0000 & 8.33 & 8.33 \\
052 & 37387 & 37387 & 37387 & 1.0000 & 1.0000 & 1.0000 & 26.03 & 26.03 \\
053 & 492232 & 492232 & 491765 & 1.0000 & 1.0009 & 1.0009 & 35.50 & 43.81 \\
054 & 516442 & 517960 & 513641 & 0.9971 & 1.0084 & 1.0055 & 0.00 & 0.00 \\
055 & 25635 & 25635 & 25635 & 1.0000 & 1.0000 & 1.0000 & 22.16 & 22.16 \\
056 & 265710 & 265710 & 265710 & 1.0000 & 1.0000 & 1.0000 & 0.00 & 0.00 \\
057 & 14028 & 14028 & 14028 & 1.0000 & 1.0000 & 1.0000 & 3.94 & 3.94 \\
058 & 1004819 & 1004819 & 1004819 & 1.0000 & 1.0000 & 1.0000 & 0.00 & 0.00 \\
059 & 150782 & 150736 & 150736 & 1.0003 & 1.0000 & 1.0003 & 73.79 & 73.79 \\
060 & 175702 & 175626 & 175626 & 1.0004 & 1.0000 & 1.0004 & 52.58 & 52.58 \\
061 & 22283 & 22283 & 22283 & 1.0000 & 1.0000 & 1.0000 & 12.21 & 12.21 \\
062 & 719016 & 718939 & 701554 & 1.0001 & 1.0248 & 1.0249 & 0.05 & 6.16 \\
063 & 251581 & 251581 & 251581 & 1.0000 & 1.0000 & 1.0000 & 0.00 & 0.00 \\
064 & 19116 & 19116 & 19116 & 1.0000 & 1.0000 & 1.0000 & 0.00 & 0.00 \\
065 & 11396 & 11396 & 11396 & 1.0000 & 1.0000 & 1.0000 & 6.64 & 6.64 \\
066 & 237259 & 237121 & 237121 & 1.0006 & 1.0000 & 1.0006 & 26.18 & 26.18 \\
067 & 15931687 & 15931687 & 15229262 & 1.0000 & 1.0461 & 1.0461 & 0.00 & 9.63 \\
068 & 341479 & 339915 & 339915 & 1.0046 & 1.0000 & 1.0046 & 9.64 & 9.64 \\
069 & 13375 & 13355 & 13239 & 1.0015 & 1.0088 & 1.0103 & 31.09 & 28.44 \\
070 & 26275 & 26216 & 26030 & 1.0023 & 1.0071 & 1.0094 & 28.70 & 29.71 \\
071 & 17557 & 17557 & 17401 & 1.0000 & 1.0090 & 1.0090 & 0.00 & 0.00 \\
072 & 4831281 & 4809827 & 4809827 & 1.0045 & 1.0000 & 1.0045 & 20.41 & 20.41 \\
073 & 3517515 & 3517515 & 3434743 & 1.0000 & 1.0241 & 1.0241 & 0.00 & 11.05 \\
074 & 337275 & 335626 & 335626 & 1.0049 & 1.0000 & 1.0049 & 31.17 & 31.17 \\
075 & 1219033 & 1219026 & 1219026 & 1.0000 & 1.0000 & 1.0000 & 17.29 & 17.29 \\
076 & 1239374 & 1232497 & 1232497 & 1.0056 & 1.0000 & 1.0056 & 39.71 & 39.71 \\
077 & 115617 & 115617 & 115617 & 1.0000 & 1.0000 & 1.0000 & 7.44 & 7.44 \\
078 & 663941 & 578013 & 578013 & 1.1487 & 1.0000 & 1.1487 & 34.36 & 34.36 \\
079 & 6075976 & 6075472 & 6075472 & 1.0001 & 1.0000 & 1.0001 & 24.42 & 24.42 \\
080 & 113455 & 90408 & 88140 & 1.2549 & 1.0257 & 1.2872 & 34.41 & 36.40 \\
081 & 30943 & 30943 & 30943 & 1.0000 & 1.0000 & 1.0000 & 0.00 & 0.00 \\
082 & 596513 & 596237 & 592853 & 1.0005 & 1.0057 & 1.0062 & 15.20 & 15.21 \\
083 & 442480 & 310731 & 310731 & 1.4240 & 1.0000 & 1.4240 & 67.57 & 67.57 \\
084 & 503472 & 503472 & 503472 & 1.0000 & 1.0000 & 1.0000 & 0.00 & 0.00 \\
085 & 2791092 & 3157239 & 3157239 & 0.8840 & 1.0000 & 0.8840 & 1.25 & 1.25 \\
086 & 118403 & 118403 & 115496 & 1.0000 & 1.0252 & 1.0252 & 0.00 & 0.00 \\
087 & 833349 & 829669 & 829449 & 1.0044 & 1.0003 & 1.0047 & 30.22 & 30.39 \\
088 & 196615 & 196344 & 194762 & 1.0014 & 1.0081 & 1.0095 & 16.01 & 15.31 \\
089 & 147432 & 147432 & 144388 & 1.0000 & 1.0211 & 1.0211 & 0.00 & 16.23 \\
090 & 542638 & 509924 & 416535 & 1.0642 & 1.2242 & 1.3027 & 8.14 & 31.93 \\
091 & 1882479 & 1876653 & 1876653 & 1.0031 & 1.0000 & 1.0031 & 31.18 & 31.18 \\
092 & 1649430 & 1351088 & 1351088 & 1.2208 & 1.0000 & 1.2208 & 58.07 & 58.07 \\
093 & 15876 & 15876 & 15876 & 1.0000 & 1.0000 & 1.0000 & 0.00 & 0.00 \\
094 & 21573 & 20380 & 20380 & 1.0585 & 1.0000 & 1.0585 & 30.61 & 30.61 \\
095 & 420852 & 420852 & 420852 & 1.0000 & 1.0000 & 1.0000 & 0.00 & 0.00 \\
096 & 41196 & 41196 & 38946 & 1.0000 & 1.0578 & 1.0578 & 20.82 & 6.94 \\
097 & 2102348 & 2102250 & 2102250 & 1.0000 & 1.0000 & 1.0000 & 24.52 & 24.52 \\
098 & 89756 & 89756 & 89756 & 1.0000 & 1.0000 & 1.0000 & 0.00 & 0.00 \\
099 & 336826 & 327731 & 327139 & 1.0278 & 1.0018 & 1.0296 & 52.40 & 52.19 \\
100 & 87316 & 87316 & 59623 & 1.0000 & 1.4645 & 1.4645 & 2.74 & 27.50 \\
\end{longtable}

\endgroup

\section{Cache参数敏感性}\label{app:sensitivity}
六图各在四个非默认硬件配置下固定五核问题二计划，并尝试候选重优化。
下表的“失败候选”表示未通过依赖检查的记录；这些记录不赋予工期。
\begingroup
\small
\setlength{\tabcolsep}{4.5pt}
\renewcommand{\arraystretch}{0.94}
\setlength{\LTleft}{\fill}\setlength{\LTright}{\fill}
\begin{longtable}{crrrrr}
\caption{6张代表图的Cache参数敏感性}\label{tab:sensitivity-detail}\\
\toprule
图 & 参数 & 固定计划/周期 & 重优化/周期 & 固定命中/\% & 失败候选 \\
\midrule
\endfirsthead
\multicolumn{6}{c}{续表}\\
\toprule
图 & 参数 & 固定计划/周期 & 重优化/周期 & 固定命中/\% & 失败候选 \\
\midrule
\endhead
\midrule\multicolumn{6}{r}{续下页}\\
\endfoot
\bottomrule
\endlastfoot
012 & 带宽×2 & 13406 & 13406 & 7.68 & 0 \\
012 & 带宽/2 & 13406 & 13406 & 7.68 & 0 \\
012 & 容量×2 & 13406 & 13406 & 7.68 & 0 \\
012 & 容量/2 & 13406 & 13406 & 7.68 & 0 \\
016 & 带宽×2 & 7706495 & 7706472 & 12.62 & 0 \\
016 & 带宽/2 & 7706495 & 7706472 & 12.62 & 0 \\
016 & 容量×2 & 7706495 & 7706472 & 24.27 & 0 \\
016 & 容量/2 & 7706495 & 7706472 & 4.85 & 0 \\
051 & 带宽×2 & 541017 & 541017 & 8.33 & 0 \\
051 & 带宽/2 & 541017 & 541017 & 8.33 & 0 \\
051 & 容量×2 & 541017 & 541017 & 16.67 & 0 \\
051 & 容量/2 & 541017 & 541017 & 0.00 & 0 \\
062 & 带宽×2 & 718939 & 701554 & 0.05 & 0 \\
062 & 带宽/2 & 718939 & 701555 & 0.05 & 0 \\
062 & 容量×2 & 718939 & 701554 & 0.05 & 0 \\
062 & 容量/2 & 718939 & 701554 & 0.05 & 0 \\
067 & 带宽×2 & 15931687 & 15218421 & 0.00 & 0 \\
067 & 带宽/2 & 15931687 & 15285122 & 0.00 & 0 \\
067 & 容量×2 & 15931687 & 15061291 & 0.00 & 0 \\
067 & 容量/2 & 15931687 & 15363155 & 0.00 & 0 \\
081 & 带宽×2 & 30943 & 30943 & 0.00 & 0 \\
081 & 带宽/2 & 30943 & 30943 & 0.00 & 0 \\
081 & 容量×2 & 30943 & 30943 & 0.00 & 0 \\
081 & 容量/2 & 30943 & 30943 & 0.00 & 0 \\
\end{longtable}

\endgroup

\section{计算参数与配对字段}
正式评价使用题面默认的L1、UB、DDR和L2 Cache参数；敏感性实验仅改变L2容量或读带宽，其他计划和硬件条件保持不变。C初始命中率与C最终命中率分别保存，固定计划和重调度计划的字段由此分开。未通过依赖检查的候选不进入工期统计。

